% ---------------------------------------------------------------------
%  TFPT --- The Prime Front (research documentation)
%  The complete prime / zeta line of the TFPT project, A to Z: the seven
%  load-bearing modules (v535-v541), the sandbox localization program,
%  the handover induction, the certification arc, the kill list, and the
%  remaining residual.  This is a RESEARCH DOCUMENTATION, not a results
%  paper: it separates what is machine-verified and ledger-carried from
%  what is exploratory sandbox work, and it claims NO progress on the
%  Riemann Hypothesis.
% ---------------------------------------------------------------------
\documentclass[11pt,a4paper]{article}
\usepackage[T1]{fontenc}
\usepackage[utf8]{inputenc}
\usepackage[english]{babel}
\usepackage{lmodern}
\usepackage[a4paper,margin=2.3cm]{geometry}
\usepackage{graphicx}
\usepackage{amsmath,amssymb,amsthm}
\usepackage{mathtools}
\usepackage{booktabs,array,tabularx}
\usepackage{xltabular}
\usepackage{enumitem}
\usepackage{xcolor}
\usepackage[most]{tcolorbox}
\usepackage{microtype}
\usepackage[hidelinks,colorlinks=true,linkcolor=blue!55!black,urlcolor=blue!55!black]{hyperref}
\emergencystretch=3em

\input{tex-artefacts/tfpt_style}
\usetikzlibrary{decorations.pathreplacing}
\newcommand{\cthree}{c_3}
\newcommand{\gcar}{g_{\mathrm{car}}}
\newcommand{\Z}{\mathbb Z}
\newcommand{\R}{\mathbb R}
\newcommand{\PP}{\mathbb P}
\newcommand{\Qcert}{Q_{\mathrm{cert}}}
\newcommand{\Afam}{A_{\mathrm{fam}}}
\newcommand{\Ashift}{A_{\mathrm{shift}}}
\newcommand{\Darch}{\Delta_{\mathrm{arch}}}
\newcommand{\Qweil}{Q_{\mathrm{Weil}}}
\newcommand{\Qfull}{Q_{\mathrm{full}}}
\setlist{itemsep=2pt,topsep=3pt}
\newcolumntype{Y}{>{\raggedright\arraybackslash}X}

% sandbox / verified provenance tags used throughout the body
\newcommand{\sbx}{{\small\textcolor{tfgray}{\textsf{[sandbox]}}}}
\newcommand{\ver}{{\small\textcolor{tfgreen!60!black}{\textsf{[ledger]}}}}
\newcommand{\schem}{\textit{Schematic.}\ }

\title{\vspace{-1.2cm}\textbf{TFPT --- The Prime Front}\\[2mm]
\large Research documentation of the prime / zeta line: seven verified modules,
one localized residual, and the complete kill list}
\author{Stefan Hamann}\date{\TFPTdatestamp}

\def\TFPTdoctitle{The Prime Front}
\begin{document}
\maketitle

\input{tex-artefacts/tfpt_docset}
\noindent\emph{You are reading the \textbf{Prime Front} companion (P): the complete internal record
of the prime / zeta line --- how a discrete compiler's bookkeeping kept producing classical
number-theoretic objects, which of those links became load-bearing modules, and how the remaining
question was compressed, over a hundred exploratory probes, down to one strict-positivity margin
input --- then, once that margin was measured to be an artifact of the requirement, down to a
two-object core --- and finally down to a three-point list of which only one point is an
inequality. It is a
\textbf{research documentation}, not a results paper. Its purpose is that a reader (or a reviewer)
can see the whole route, including everything that died on the way.}

\begin{abstract}
\noindent A two-axiom discrete compiler ($\cthree=1/(8\pi)$, $\gcar=5$, $D_5\oplus A_3+\mu_4\Rightarrow
E_8$) has no business meeting the prime numbers. It does anyway: the $\mu_4$ glue completion carries a
quartic character, the signed glue census is a theta series, and the shell counts of that census turn
out to be \emph{Hecke} data. It goes further than an analogy: the $E_8$ shell count is literally
$240\,\zeta(s)\zeta(s-3)$ as a Dirichlet series, and the lattice's self-duality gives its theta sum an
exact functional equation whose fixed axis is the critical line --- a framework picture, stated as such,
which the document opens with and then keeps at arm's length. This companion documents the line end to
end. Seven modules are
load-bearing and ledger-carried (\vref{v535}--\vref{v541}): Hecke from geometry on the $E_8$ census,
the Eichler trace layer, the half-integral bridge, one finite relative-trace identity of which the
first three are projections, the Weil structure of the resulting family with two explicitly isolated
obstructions, the amplitude route with a positive linear carrier, and the matching-lemma / transport
ledger package. Everything beyond those seven is \emph{sandbox}: an exploratory program, run as
$115$ preregistered probes with $2876$ machine checks, which localized the residual object. That
localization is the actual content of this document. The transport ledger closes exactly, leaving one
coupled inequality (called \textbf{I5}) which is --- \emph{by identity, as a typing and not as
progress} --- equivalent to Weil positivity and hence to the Riemann Hypothesis. The subsequent arc
maps I5 geographically (Connes--Consani dictionary, tightness curves, the Sonin crossing, the one-atom
band), turns the positivity question into a per-zone handover induction, and then certifies its
ingredients one by one until the entire remaining hardness sits in a single Loewner statement on the
$J$-odd half of a window, reducible by Sherman--Morrison/Woodbury to one scalar ratio
$r=\kappa/\varepsilon\le1$ (measured $0.0053$--$0.1814$ on $16$ zones), then, once $\varepsilon$ is
identified as an exact energy, to two scalars --- and finally, with both scalars certified (the wing
scalar unconditionally by a graded matrix cap, the boundary value by a residual certificate that
carries a cancellation), to a chain that closes $16/16$ conditional on exactly one strict-positivity
margin input, itself two to six orders of magnitude weaker than the conclusion it buys. The
penultimate part then closes the induction circle on the measured zones end to end --- a certified
base case, fifteen certified handover steps through a graded Loewner minorant, and the structural
finding that the atom entry costs nothing --- while leaving three sharp, named gaps: no reserve, no
scalar step law, and no uniformity in $k$, with an extrapolated crossing warning at $n\approx170$.
The next part stops extrapolating and measures: on a deep ladder of $199$ prime-power zones with
$117$ certified handovers, the crossing is located between $n=461$ and $n=463$ --- the $n\approx170$
figure was a fit artefact, but the wall is real, and $n^*\approx462$ is an upper bound --- and the
single warning splits into three separate walls (a measured margin wall, an arithmetic ladder wall
set by the twin-prime pair $521\to523$, and a vacuous-requirement wall at $n=727$), while the
handover mechanism itself never fails: all $117$ steps certify at retention $1.000000$, and the
chain tears at the ratio, not at a step. The operating variable turns out to be the window depth,
not the zone index. The next part acts on that reinterpretation and rebuilds the construction in a
gap-coupled scaled frame ($D$ tied to the local prime gap): two of the three walls fall
\emph{structurally} --- the ladder death becomes a payable cell cost (the killer pair $461\to463$
now certifies spectrally with the reserve opening by a factor ${\sim}1000$) and the requirement
wall loses its depth dependence --- while the margin wall proves \emph{frame-invariant} at exponent
$-0.974$: it is the substance of the requirement, not geometry. The hardness is thereby traded ---
three arithmetic walls against two analytic demands, the positivity of one limit operator and a
convergence rate below the step reserve --- with the prime-gap dependence disclosed rather than
hidden. The next part settles the currency question, with a fundamental reinterpretation: the wall
carries the same exponent in \emph{every} admissible currency --- it is substance --- but it
measures the \emph{discretisation}, not the spectrum. Under refinement at a fixed window both
leading eigenvalues carry the same cell-width power (the continuum window form has no gap), the
positive floor survives only as a cancellation of relative size $10^{-7}$--$10^{-4}$, and the T109
requirement chain turns out to divide by an artifact margin; the hardness balance is restated as
four parts, led by a margin-free step certificate. The next part builds that certificate and the
wall dissolves: the exact Schur complement (Albert 1969, with the Moore--Penrose inverse)
certifies every ladder step margin-free --- eleven of them beyond the old wall, to $n=1331$, and
all seven zones where the requirement chain tore --- the wall was an $O(1)$ numerator divided by
an artifact floor, chains now stop only at the compute cap and never at a step, and the (R)
demand itself proves grid-bound and is deleted rather than transported, leaving a two-object
core: $\varepsilon$ relative to $\kappa$ via the exact Szeg\H{o} identity, and transport of an
$O(0.1)$ Schur complement between grids. The last part splits that core: the transport across a
real regrid certifies only mild refinement (a clean split at ratio $\rho^*\approx1.8$, and for a
principled reason --- on nested ladders, where the transport error is exactly zero, the Schur
floor itself falls like $\rho^{-1.7}$, so no bound can undo the drop), while a two-scale
compression breaks the compute cap: a certified margin-free step at $n=155\,921$ ($117\times$
deeper), chains of ten steps, the stopper always cost and never a step --- leaving the shortest
list the program has produced: three points, of which only one is an inequality (the classical
Szeg\H{o}--Levinson prediction-error bound for one symbol). \textbf{No claim of progress on the Riemann Hypothesis is
made anywhere in this document}, and none of the sandbox findings moves any ledger marker. The
document is deliberately symmetric about failure: a dedicated section lists the readings and
routes that were refuted --- an essential-singularity reading, a $C/g$ proxy law, a density chain, an
accumulated invariant with self-propagation, the Landau--Widom anchor, the symbol route for the odd
channel, three certificate candidates for the avoidance norm plus the Combes--Thomas certificate at
the boundary, and the naive margin route --- because in
a program of this shape the kill list is the more reliable output.
\end{abstract}

\TFPTbadges

% =====================================================================
\section{How to read this document}
% =====================================================================
\begin{gapboxnb}[Scope and status --- read this before anything else]
\begin{enumerate}[leftmargin=1.4em,label=\textbf{\arabic*.}]
\item \textbf{No RH claim.} Nothing in this document claims progress towards, or evidence for, the
Riemann Hypothesis. Where an equivalence to RH appears (Section~\ref{sec:i5}) it is an
\emph{equivalence typing}: a statement about which object carries the difficulty, produced by an exact
identity, and deliberately \emph{not} a statement that the difficulty has been reduced.
\item \textbf{Two layers, never mixed.} Claims marked \ver\ are load-bearing: they have a row in
\texttt{verification/status\_ledger.csv}, a module in \texttt{verification/}, and a script citation
here (the blue \texttt{[\,verification/\dots\,]} tag). Everything marked \sbx\ is exploratory work in \texttt{experiments/tfpt-discovery/}; it carries
no ledger row, no marker, and no authority. Section~\ref{sec:verified} is the only section that is
entirely \ver.
\item \textbf{Status markers are quoted, never upgraded.} The markers \mE\ \mC\ \mO\ \mX\ shown in
Section~\ref{sec:verified} are read verbatim from the ledger rows of \vref{v535}--\vref{v541}. This
document introduces no new markers and moves none.
\item \textbf{Fits are declared as fits.} Every scaling law, slope, exponent and extrapolation quoted
from the sandbox is a fit to finitely many measured windows and is labelled as such at the point of
use. Sixteen resolved zones are sixteen zones, not all zones.
\item \textbf{Classical results are named as classical.} Kneser neighbours, Eichler/Witt, Shimura,
Waldspurger, Cohen, Weil/Guinand, Bombieri, Yoshida, Connes--Consani, Gronwall/Robin, Douglas,
Slepian--Landau--Pollak, Sherman--Morrison/Woodbury, Grenander--Szeg\H{o} and the rest are classical
and are used as such; the compiler-native part is the wiring, not the theorems.
\item \textbf{Numbers are literal.} Every numeric value in this document is transcribed from the
research diary \texttt{experiments/next.txt} or from a ledger row. Figures that could not be drawn
from literal data are marked \emph{schematic} in their caption, and the caption states exactly which
quantities in the figure are literal.
\end{enumerate}
\end{gapboxnb}

\noindent The diary parts are cited as \textbf{T$n$} (part $n$ of the research diary, chronological,
$n=11\ldots115$). The probe file, the contract name, the verdict and the check count of every part are
listed in Appendix~\ref{app:index}. Section~\ref{sec:motivation} explains why a physics compiler meets
the primes at all and Section~\ref{sec:bigpicture} states the framework interpretation that motivated
the whole line --- a self-dual lattice, its functional equation, and a mirror that returns at the far
end of the program; Section~\ref{sec:verified} is the verified layer; Sections~\ref{sec:i5},
\ref{sec:handover} and \ref{sec:certification} are the sandbox program in three acts;
Section~\ref{sec:kills} is the kill list; Section~\ref{sec:status} is the status statement.

% =====================================================================
\section{Motivation: why a compiler meets the primes}\label{sec:motivation}
% =====================================================================
TFPT is a discrete compiler with two inputs: the seam constant $\cthree=\tfrac1{8\pi}$ (P1) and the
carrier rank $\gcar=5$ (P2). They force the lattice completion $D_5\oplus A_3+\mu_4\Rightarrow E_8$,
and the readouts of the Standard-Model skeleton are taken off that completion. None of this is about
number theory. The prime line begins with an accident of the \emph{glue}.

The completion is not $D_5\oplus A_3$ but $D_5\oplus A_3$ \emph{plus} a glue group $\mu_4$ --- a cyclic
group of order four. A cyclic group of order four carries a quartic character, and the quartic
character on $\Z$ is $\chi_{-4}$, the non-trivial character mod $4$. Once the census of glue vectors is
counted \emph{with that sign}, three things happen at once, and all three are elementary consequences
rather than choices:

\begin{enumerate}[leftmargin=1.5em]
\item The signed glue count is a \textbf{theta series}. Signed counting of a quadratic form is a theta
series by definition; here the compiler's own count $\theta_3^2$ is exactly the Gaussian-integer theta,
$r_2(n)$, so the census counts $\Z[i]$-ideals (T85).
\item The character makes the count a \textbf{Dirichlet object}. $\chi_{-4}$ is the character of
$\Q(i)$; the associated $L$-function is the first Dirichlet $L$-function, and the split law
``$p=a^2+b^2$'' becomes a compiler-visible fibre property of each prime (T21, channel B: $\chi_4$ fibre
$\iff p=a^2+b^2$, verified $81/0/87/0$ on $p<1000$; $\chi_3$ fibre $\iff p=a^2+3b^2$, $80/0/86/0$).
\item The shell counts of the census carry \textbf{Hecke eigenvalues}. That is the surprise, and it is
the content of the first load-bearing module.
\end{enumerate}

\begin{figure}[ht]\centering
\begin{tikzpicture}[
  font=\small,
  box/.style={draw=tfblue!70!black,fill=tfblue!5,rounded corners=2pt,inner sep=4pt,align=center,
              text width=31mm,minimum height=11mm},
  gbox/.style={box,draw=tfgreen!65!black,fill=tfgreen!6},
  ar/.style={-{Latex[length=2mm]},draw=tfgray!80,thick}]
\node[box]  (ax)      at (0,0)     {two axioms\\[-1pt]\footnotesize $\cthree=\tfrac1{8\pi}$, \ $\gcar=5$};
\node[box]  (e8)      at (3.9,0)   {$D_5\oplus A_3+\mu_4$\\[-1pt]$\Rightarrow E_8$};
\node[box]  (glue)    at (7.8,0)   {glue group $\mu_4$\\[-1pt]\footnotesize quartic character $\chi_4$};
\node[box]  (zi)      at (11.7,0)  {$\Z[i]$, \ $\theta_3^2=r_2$\\[-1pt]\footnotesize signed census};
\node[gbox] (weight4) at (3.9,-2.6){weight-4 newform $f_8$\\[-1pt]\footnotesize $\eta(2\tau)^4\eta(4\tau)^4$};
\node[gbox] (hecke)   at (7.8,-2.6){Hecke data on the census\\[-1pt]\footnotesize $\nu_p=a\,\mathrm{Id}+b\,T_p$};
\node[box]  (lfun)    at (11.7,-2.6){Dirichlet $L(s,\chi_{-4})$\\[-1pt]\footnotesize split law $p=a^2+b^2$};
\draw[ar] (ax)    -- (e8);
\draw[ar] (e8)    -- (glue);
\draw[ar] (glue)  -- (zi);
\draw[ar] (glue)  -- (hecke);
\draw[ar] (zi)    -- (lfun);
\draw[ar] (hecke) -- (weight4);
\draw[ar] (hecke) -- (lfun);
\end{tikzpicture}
\caption{The chain by which the compiler acquires arithmetic. Nothing here is chosen for number
theory: the glue group is forced by the completion, the quartic character is the character of a cyclic
group of order four, and the signed count of a quadratic form is a theta series. The green nodes are
where the line becomes load-bearing (Section~\ref{sec:verified}).}\label{fig:chain}
\end{figure}

\noindent The first question of the series was therefore not ``does this prove anything about
primes'' but ``which of these links are \emph{mechanism} and which are pretty coincidences''. Parts
T11--T13 established the signed glue theta, its $L$-series and the glue-character atlas
($19/19$, $16/16$, $11/11$); T14--T19 tested six candidate contracts in parallel and produced two
self-corrections against the project's own earlier readings (T14 \textsc{deflated}: the measured
$2\pi$ is the universal Bisognano--Wichmann/Unruh conversion, not ``the self-dual width''). The
cross-finding of that batch is worth stating because it survived the entire series: the structure is
exact at the census/lattice level (weight $4$), and the gap is always the \emph{same} gap --- the
canonical transport to the $\mathrm{GL}(1)$ / $\xi(s)$ world (multiplicity one, weight $1/2$, the
archimedean place, the derivation of the Hecke action). One meta-contract, \texttt{ZETA.CENSUS.TO.GL1}
\mO.

Parts T20--T25 then killed, in four stages, the slogan ``the archimedean term comes from the seam'':
the raw seam density of states is flat/slightly decreasing (loglog slope $\approx-0.40$) while the
classical digamma kernel grows logarithmically (measured $0.15917$ against $1/(2\pi)=0.159155$, rms
$3\times10^{-6}$ for $t\ge20$) --- \textsc{killed-as-naive} (T20, $25/25$); the non-DOS successor
located the missing logarithm in the interval cut but failed the one-constant dictionary at $2/\pi$
(T22, \textsc{partial}); and both dictionary closures plus the scattering-phase observable died
(T23--T25, $23/23+25/25+29/29$). This matters for the rest of the document: the archimedean place was
never going to come for free, and when it finally arrived (T82) it arrived from \emph{inside} the same
family, not from the seam.

% =====================================================================
\section{The big picture: a self-dual lattice and its own mirror}\label{sec:bigpicture}
% =====================================================================
\begin{gapboxnb}[What this section is]
A \textbf{framework interpretation}, not a result. It says how the pieces fit together and why the line
was pursued for a hundred parts. Every factual statement it uses is either classical (and named as
such), one of the seven modules of Section~\ref{sec:verified}, or a sandbox probe with its part number.
Nothing here is an argument about the Riemann Hypothesis, nothing here is used later as a premise, and
the picture itself carries no status marker.
\end{gapboxnb}

\subsection{(i) The zeta function is literally in the geometry}
Count the $E_8$ lattice points shell by shell. The even shells satisfy $N(2n)=240\,\sigma_3(n)$ ---
direct lattice counting gives $240$, $2160$, $6720$, $17\,520$, $30\,240$, $60\,480$ for $n=1\ldots6$
(T6, one of the parts that predate the indexed range of Appendix~\ref{app:index}) --- and $\sigma_3$ is
multiplicative, so the counting function of the lattice is a Dirichlet series
with an Euler product:

\begin{keyeq}$\displaystyle
\sum_{n\ge1}\frac{N(2n)}{n^{s}}\;=\;240\sum_{n\ge1}\frac{\sigma_3(n)}{n^{s}}\;=\;240\,\zeta(s)\,\zeta(s-3)
\qquad(\text{at }s=5:\ 1.705672\ \text{vs}\ 1.705678).$\end{keyeq}

\noindent The Riemann zeta function stands in $E_8$'s own counting function, with a simple pole at
$s=4$ of residue $240\,\zeta(4)=8\pi^4/3$ (T12). This is classical lattice theory --- Eisenstein series
of weight $4$ --- and it is stated here only to fix where the arithmetic enters: not through a modelling
choice, but through the object the compiler already had. The contrast is instructive: the same
multiplicativity \emph{fails} for the $D_5\oplus A_3$ sublattice alone ($52$, $576$, $1584$, $16\,128$
with $576\cdot1584\neq16\,128\cdot52$, T6). Without the glue there is no Euler structure in the count.

\subsection{(ii) The glue carries the character of the Gaussian integers}
The completion is $D_5\oplus A_3+\mu_4$, and a cyclic group of order four carries a quartic character,
which on $\Z$ is $\chi_{-4}$ --- the character of $\Q(i)$. Two consequences make the primes a
\emph{fibre property} of the compiler rather than an application of it. First, the split
$p\equiv1$ versus $p\equiv3\pmod4$ is visible to the census directly: the $\chi_4$ fibre holds exactly
for $p=a^2+b^2$ ($81/0/87/0$ on $p<1000$, T21). Second, the character kills a pole. Of the three
character channels of the $E_8$ counting function, the trivial one is $240\zeta(s)\zeta(s-3)$ with its
pole at $s=4$, the spinor one is $64(1-2^{-s})\zeta\zeta$ and keeps it, while the $\mu_4$-signed channel
$16\cdot2^{-s}\eta(s)\eta(s-3)-8L(f_8,s)$ is \emph{entire} ($L_c(4.001)=-6.976$, finite; T12). The
signed channel is the only pole-free one --- exactly the $\zeta$ versus $L(\chi)$ pattern.

\subsection{(iii) Hecke structure is an output of geometry, not an input}
The step that made the line load-bearing is \vref{v535}: the Hecke action on the census is produced by
\emph{pure lattice adjacency} --- Kneser $p$-neighbours --- with no modular input anywhere in the
derivation. The marked neighbour sum is exactly an affine Hecke element, $\nu_p=a\,\mathrm{Id}+b\,T_p$,
and the cusp eigenvalues come out of the glue census as $a_p=-c(p)/8$. Prime structure is what the
geometry \emph{emits}; nothing about primes was put in. The first four modules of
Section~\ref{sec:verified} make this precise and end in one finite relative-trace identity of which the
first three are projections.

\subsection{(iv) The core point: self-duality, and a mirror that keeps coming back}
$E_8$ is unimodular --- self-dual --- and self-duality is exactly the hypothesis of Poisson summation.
The lattice theta sum therefore satisfies an exact functional equation,
$\sum_x e^{-t|x|^2}=(\pi/t)^4\sum_x e^{-(\pi^2/t)|x|^2}$ (verified to $25$ digits, T12), whose unique
self-dual width is $t=\pi$. A functional equation is a \emph{reflection}, and the critical line is the
fixed axis of that reflection. That is the geometric content of the picture: on the axis sits the
symmetry, and beside it sits everything the lattice actually carries.

Two findings, at opposite ends of the program, put the same mirror on the table twice.

\begin{itemize}[leftmargin=1.5em]
\item \textbf{The reflection product is the transport operator} (T83,
\texttt{fe\_symmetrization\_probe.py}, $27/27$, \textsc{invariant-null}). The involution algebra was
computed exactly: $\zeta$'s functional equation is the reflection $J_{1/2}$ and is absorbed in evenness
(on autocorrelations it acts as the identity), while the family's functional equation is $J_1$. Their
\emph{product} $J_1\circ J_{1/2}=e^{\pm u}$ is the unit-line shift --- the difference of the two centres
\emph{is} the transport operator, algebraically. The two reflections generate an infinite dihedral
ladder; the value side is invariant under the whole ladder, the spectral cone only under its own
reflection, which is what anchors it to the critical line; and $\mathrm{Fix}(J_1)\cap\{\text{Weil
cone}\}=\{0\}$ exactly, so no symmetrisation can relieve I5. The verdict was a null with a structural
finding: the wall is \emph{transversal} to the functional equation, not positional.
\item \textbf{The same mirror sits in the last hard residue} (T106, Section~\ref{sec:certification}).
The mirror parity $J$ with $JQJ=Q$ and $JB=-BR$ splits the Weil pole exactly into
$ss^{\mathsf T}-tt^{\mathsf T}$: a positive rank-one lift in the $J$-even channel and a negative
rank-one pressure in the $J$-odd one. The even channel closes; everything that is still hard lives in
the odd half --- the anti-symmetric part of the same reflection.
\end{itemize}

\noindent Stated as a picture and nothing more: the recurring question of this line is whether the
self-dual object is entirely at peace with its own mirror. That sentence is an interpretation of where
the difficulty sits, in the same spirit as the classical reading of the critical line as a symmetry
axis. It is \emph{not} a reformulation of the Riemann Hypothesis, it is not used as a premise anywhere
below, and it produces no claim.

\subsection{(v) What the program did with the picture}
It converted it, step by step, into smaller and more explicit objects. The transport ledger closes
exactly (\vref{v541}), which leaves the single coupled inequality I5; I5 is by identity equivalent to
Weil positivity, an equivalence \emph{typing} and not a reduction (Section~\ref{sec:i5}); the positivity
question is then reorganised as a per-zone handover induction (Section~\ref{sec:handover}); the
induction is certified ingredient by ingredient until only one Loewner statement (R) on the $J$-odd half
of a window remains; and (R) reduces first to one scalar ratio, then to two scalars, and finally ---
with both scalars certified --- to one strict-positivity margin input
(Section~\ref{sec:certification}). The full
cascade, with the type of every arrow marked, is Figure~\ref{fig:cascade}; the reason it is a route and
not a proof sketch is stated there and in Section~\ref{sec:status}.

\begin{figure}[ht]\centering
\begin{tikzpicture}[font=\small,
  ax/.style={draw=tfgold!75!black,dashed,very thick},
  sb/.style={draw=tfblue!70!black,fill=tfblue!5,rounded corners=2pt,inner sep=5pt,align=center,
             text width=52mm,minimum height=20mm},
  fb/.style={draw=tfgreen!65!black,fill=tfgreen!6,rounded corners=2pt,inner sep=5pt,align=center,
             text width=52mm,minimum height=22mm},
  qb/.style={draw=tfgray!70,fill=tfgray!6,rounded corners=2pt,inner sep=5pt,align=center,
             text width=118mm,minimum height=10mm},
  ar/.style={-{Latex[length=2.2mm]},draw=tfgray!85,thick},
  bi/.style={{Latex[length=2.2mm]}-{Latex[length=2.2mm]},draw=tfgold!75!black,thick}]
\draw[ax] (0,1.45) -- (0,-5.45);
\node[align=center,font=\footnotesize\itshape,text=tfgold!55!black,text width=42mm] at (0,2.15)
  {the fixed axis of the self-duality};
\node[sb] (l) at (-3.65,0.1) {\textbf{one side}\\[2pt]\footnotesize
  $\sum_x e^{-t|x|^2}$ \ at width $t$};
\node[sb] (r) at (3.65,0.1) {\textbf{the mirrored side}\\[2pt]\footnotesize
  $(\pi/t)^4\sum_x e^{-(\pi^2/t)|x|^2}$};
\draw[bi] (l) -- (r);
\node[fill=white,inner sep=2pt,font=\footnotesize,text=tfgold!55!black] at (0,0.1) {$J$};
\node[fb] (t83) at (-3.65,-2.85) {\textbf{T83} --- the mirror as algebra\\[2pt]\footnotesize
  $J_1\circ J_{1/2}=e^{\pm u}$ is the unit-line shift: the reflection product \emph{is} the
  transport operator};
\node[fb] (t106) at (3.65,-2.85) {\textbf{T106} --- the mirror as obstruction\\[2pt]\footnotesize
  the pole splits $ss^{\mathsf T}-tt^{\mathsf T}$; the even channel closes, the odd one does not};
\node[qb] (q) at (0,-5.9) {\footnotesize\emph{Framework interpretation:} the recurring question is
  whether the self-dual object is at peace with its own mirror. No claim.};
\draw[ar] (l) -- (t83);
\draw[ar] (r) -- (t106);
\draw[ar] (t83) -- (q);
\draw[ar] (t106) -- (q);
\end{tikzpicture}
\caption{\schem The self-duality picture. \emph{Literal:} the Poisson identity and its unique self-dual
width $t=\pi$ ($25$ digits, T12); the algebraic identity $J_1\circ J_{1/2}=e^{\pm u}$ and
$\mathrm{Fix}(J_1)\cap\{\text{Weil cone}\}=\{0\}$ (T83); the exact pole splitting
$ss^{\mathsf T}-tt^{\mathsf T}$ at $10^{-18}$ (T106). \emph{Schematic:} the geometric arrangement, and
the reading of the axis as ``order'' with the lattice's arithmetic content beside it. The bottom line is
an interpretation and carries no status.}\label{fig:bigpicture}
\end{figure}

% =====================================================================
\section{The verified layer: seven modules}\label{sec:verified}
% =====================================================================
\noindent This section is entirely \ver. Each subsection states the claim as the ledger states it,
cites the module, and quotes the ledger's own status markers. Nothing here depends on
Sections~\ref{sec:i5}--\ref{sec:certification}.

\begin{keybox}[The seven modules and their ledger markers]
\centering\small
\begin{tabularx}{\linewidth}{@{}l l Y l@{}}
\toprule
\textbf{Module} & \textbf{Claim id} & \textbf{Content} & \textbf{Markers}\\
\midrule
\vref{v535} & \texttt{HECKE.GEOM.01} & Hecke from geometry on the $E_8$ census channels & \mE\ + \mC\ + \mO\\
\vref{v536} & \texttt{HECKE.GEOM.EICHLER.01} & Eichler trace layer at the $E_8$ lattice & \mE\ + \mC\ + \mO\\
\vref{v537} & \texttt{HECKE.GEOM.HALFINT.01} & half-integral bridge & \mE\ + \mC\\
\vref{v538} & \texttt{HECKE.GEOM.RTF.01} & compiler relative trace identity & \mE\ + \mC\\
\vref{v539} & \texttt{RTF.GNS.WEIL.01} & Weil structure of the compiler family & \mE\ + \mC\\
\vref{v540} & \texttt{RTF.GNS.AMP.01} & amplitude route and positive linear carrier & \mE\ + \mC\\
\vref{v541} & \texttt{RTF.GNS.LEDGER.01} & matching lemma and transport ledger package & \mE\ + \mC\\
\bottomrule
\end{tabularx}
\markerkey
\noindent Every row carries \texttt{NO marker moves} in the ledger. The \mO\ on \vref{v535} and
\vref{v536} is the same open gate throughout: \emph{weight-4 $\to$ $\mathrm{GL}(1)$ transport}.
\end{keybox}

\subsection{Hecke from geometry (\texorpdfstring{\vref{v535}}{v535})}
The Hecke structure of the $\mu_4$-glue census is lattice-native: it comes out of Kneser
$p$-neighbour geometry, not out of modular input \veri{v535\_hecke\_from\_geometry.py}
($25$ checks, ${\sim}11$\,s; consolidated promotion of T27/T31/T32).

\begin{itemize}[leftmargin=1.5em]
\item \textbf{Kneser carries Hecke} \mE: $\#\mathrm{iso\_pts}=p^7+p^4-p^3$ and
$\#\mathrm{iso\_lines}=\sigma_3(p)\cdot\#\PP^3(\mathbb F_p)$ identically and by full enumeration at
$p=2,3,5,7$ ($135/1120/19656/137600$ lines). The Eisenstein eigenvalue $\sigma_3(p)=1+p^3$ is an exact
factor of a lattice-native correspondence count.
\item \textbf{The marked neighbour sum is an affine Hecke element} \mE:
$\nu_p=a\,\mathrm{Id}+b\,T_p$ exactly (overdetermined, residual $0$), with $b=\sigma_3+a_p$; the cusp
rule $a_p=b-\sigma_3$ recovers $a_3=-4$ and $a_5=-2$.
\item \textbf{Census redundancy is $2$-adic oldform structure} \mE: $\dim V=7=5+2$,
$\pi_{\mathrm{cusp}}=(28-T_3)/32$, only $2$-adic levels occur.
\item Fences carried \emph{inside} the claim: the $p=5$ geometric $S[1]$ block is a T31-validated
\emph{freeze} \mC; the weight-4 $\to\mathrm{GL}(1)$ transport is \mO.
\end{itemize}

\subsection{The Eichler trace layer (\texorpdfstring{\vref{v536}}{v536})}
The smooth (Eisenstein) background and the coherent (cusp) interference separate exactly
\veri{v536\_eichler\_trace\_layer.py} ($23$ checks, ${\sim}30$\,s; promotion of T33/T36/T37/T42).
The Witt layer gives $\lambda_{\mathrm{Eis}}=\sigma_3(\#\PP^3-\sigma_3)=L-\sigma_3^2$ identically
(anchors $336$, $3780$, $19264$ at $p=3,5,7$) \mE, and the Eichler identity
$\lambda_{\mathrm{geom}}=\lambda_{\mathrm{Eis}}+a_p^2$ holds two-sidedly ($352/3784/19840$). The signed
cusp reading is $a_p=-c(p)/8\Rightarrow(-4,-2,+24)$, and the assembler $R(p)=a_p^2$ is exact for
$p\le31$. The $p=5$ type-(ii) block is again a validated freeze \mC; the transport gate stays \mO.

\subsection{The half-integral bridge (\texorpdfstring{\vref{v537}}{v537})}
In the $70$-element compiler weight-$5/2$ theta monoid there is \emph{exactly one} object $g$ with
$\mathrm{Sh}_{t=2}(g)=-8f_8$ coefficientwise \mE, it is an exact $T(p^2)$-eigenform with the
weight-4 eigenvalues $(-4,-2,24,-44,22)$ \mE, and $U_4(g)=0$ places it outside the Kohnen plus space
--- a scope fence, not a defect \mE\ \veri{v537\_halfintegral\_bridge.py} ($20$ checks, ${\sim}90$\,s;
promotion of T38/T41/T44/T45). The Waldspurger/Baruch--Mao constant
$R=23.1873585645\ldots$ is constant on ten discriminants $d\equiv1\bmod 8$ (spread $\le10^{-12}$),
typed \mC\ because it is an AFE-numeric measurement.

\subsection{One finite relative trace identity (\texorpdfstring{\vref{v538}}{v538})}
The strongest structural statement of the verified layer: \vref{v535}, \vref{v536} and \vref{v537} are
three \emph{projections} of one finite relative-trace identity of Jacquet type
\veri{v538\_relative\_trace\_identity.py} ($18$ checks, ${\sim}14$\,s; promotion of T53).
$\mathrm{Tr}_V(\nu_p\circ\pi)$ is computed twice, independently --- geometrically from the glue census
with no modular input, and spectrally from $\eta(2\tau)^4\eta(4\tau)^4$ with no geometry --- and the
two agree exactly for all $(p,\pi)\in\{3,5,7\}\times\{\mathrm{Id},\pi_{\mathrm{Eis}},
\pi_{\mathrm{cusp}},\pi_{E_4},\pi_{f_8}\}$ (anchors, e.g.\ $p=7$: full trace $727680$, $f_8$ channel
$19840$) \mE. The Eichler split \emph{is} the orbit decomposition: principal orbit
$a\,\mathrm{Id}\to\lambda_{\mathrm{Eis}}$, elliptic orbit $b\,T_p\to a_p^2$, cross terms cancel exactly.
And the period side is literally compiler geometry: $g$ is the quaternary lattice form
$n=(x^2+y^2)/2+2z^2+u^2+2w^2$ with sign $(-1)^{|u|+|w|}$, giving
$[\text{signed lattice count}]^2=R\,|d|^{3/2}L(f_8\times\chi_d,2)$ at relative $10^{-12}$ \mC.
The identity is \emph{finite}; the infinite RTF is explicitly open.

\subsection{The Weil structure of the family (\texorpdfstring{\vref{v539}}{v539})}
The family carries a fully identified Weil structure \emph{up to two explicitly isolated obstructions},
and the obstructions are promoted content, not a footnote
\veri{v539\_weil\_structure\_family.py} ($25$ checks, ${\sim}6$\,s; promotion of T55/T61/T62/T63).
The canonical GNS pairing is a direct integral over the Gelfand spectrum with fibre mixing
$\le5.04\times10^{-16}$; the trivial Sato--Tate isotype is algebraically a $\mathrm{GL}(1)$ core,
$G_0=(1+Y)/(1-Y)$; the prime side is an exact linear relation to shifted $\zeta$-Weil forms at
relative $\le9\times10^{-16}$. The two obstructions are (i) a sign/doubling residue (metaplectic, $\mathrm{sym}^2$)
that survives every preregistered canonical projection, and (ii) a Plancherel correction which is
irreducibly non-automorphic. Finite-class positivity of $Q_{\mathrm{fam}}$ is measured \mC\
($[4.369,\,11.486]$); dense-class positivity is open. The ledger row says it in one line:
\emph{not almost-RH}.

\subsection{The amplitude route and the positive linear carrier (\texorpdfstring{\vref{v540}}{v540})}
The route out of the square plane \veri{v540\_amplitude\_linear\_carrier.py} ($34$ checks, ${\sim}3$\,s;
promotion of T67--T72). An amplitude Dirac operator
$D=\left(\begin{smallmatrix}0&V\\V^{\mathsf T}&0\end{smallmatrix}\right)$ with
$V[d,m]=\sqrt{w_d}\,b(d)\chi_d(m)$ exists and is Hecke-equivariant \mE; the geometric polarisation
$b=N_+-N_-$, $\Theta=N_++N_-$ is lattice-exact with $N_\pm\ge0$ for all $n\le50\,000$ \mE; the Cohen
seed $\Theta(d)=-48L(-1,\chi_d)$ is exact-rational on $159$ live $d$ \mE; the prime side is a
\emph{pure plus} combination $Q=Q_\zeta(g_-)+Q_\zeta(g_+)$ at relative $2\times10^{-15}$ \mE; the
functional equation of the linear family holds at relative $1.4\times10^{-40}$ \mE. The open boundary
is \emph{inside} the claim: the gap functional $\lambda^*$ on $n\equiv6\bmod 8$, with an exact Farkas
witness ($0.0511$), is a named limit, and positivity is asserted only in the Euler region.

\subsection{The matching lemma and the transport ledger (\texorpdfstring{\vref{v541}}{v541})}
The largest module of the line \veri{v541\_matching\_lemma\_ledger.py} ($33$ checks, ${\sim}10$\,s;
promotion of T78--T82 and T85; every core check recomputed, not cited).

\begin{itemize}[leftmargin=1.5em]
\item \textbf{Window proof} \mE: the matching lemma is proved exact-integer on $[4,10^6]$ ---
$7S(j)<40A(j)$ at every atom, $0$ violations over $939\,870$ clash atoms, no-overflow proven in exact
integers, $268$ big-integer rechecks; exact margin $X=22242575693/270725400000=0.082159$ at
$j^*=554400$; $\rho_{\mathrm{crit}}=1.143989>21/20$. A maximal-greedy identity makes the certificate
uniform in the target set and the greedy weights.\footnote{Prior-art note (literature pass,
2026-07-27): the coefficient laws L1--L4 that T77/T78 use to shape the lemma are specialisations of
Cohen 1975 (Math.\ Ann.\ 217: the Hurwitz class-number series $H(r,n)$ and the Cohen--Eisenstein
series of weight $r+\frac12$) and of Pei--Wang 2003 --- corollaries of known theory, not discoveries.
The restricted divisor inequality itself remains the independent part.}
\item \textbf{Transport ledger} \mE: $\Qweil=\Qcert+\Darch+\Delta_2$ closes with residual
$4.4\times10^{-16}$ on a $24$-row battery, and $\Delta_{\mathrm{pole}}\equiv\Delta_{\mathrm{conv}}\equiv0$
are \emph{proven}, not measured (kernel collapse in \texttt{sympy}, exact scale invariance).
\item \textbf{Character-exact envelope} \mE: $-\psi=(\chi_{-4}+\tfrac14\chi_8+\tfrac14\chi_{-8})\Theta$
coefficient-exact for all odd $n\le50\,000$; confinement holds as a \emph{set equality} on $10^6$.
\item \textbf{Archimedean term internal} \mE: the Legendre duplication bridge
$(2\pi)^{-s}\Gamma(s)=\tfrac12\Gamma_\R(s)\Gamma_\R(s{+}1)$ and the kernel identity
$k_\zeta=K_{\mathrm{fam}}-K_{\mathrm{shift}}$ at $7.9\times10^{-31}$ pointwise.
\item \textbf{Frontier facts} \mC: unlifted coherent chain crossing at $k^*=14$
($\log_{10}N^*=23.1$); Robin tail factor $6.16$ at $J=10^6$ and divergent --- the constant route stays
open.
\item \textbf{Named limits as content.} Two limits are part of the claim rather than hidden by it:
one classically-shaped open lemma (correlated cancellation on credit-rich non-coherent tail atoms), and
the I5 equivalence typing discussed next. The row states explicitly: \emph{NOT ``almost RH''};
\texttt{ZETA.HP.CARRIER} untouched.
\end{itemize}

% =====================================================================
\section{The localization program (sandbox)}\label{sec:i5}
% =====================================================================
\begin{gapboxnb}[Provenance]
Everything from here to Section~\ref{sec:status} is \sbx: exploratory probes in
\texttt{experiments/tfpt-discovery/}, no ledger row, no marker, no authority. The findings are
reported because the \emph{shape} of the residual is the useful output, not because any of them is
established.
\end{gapboxnb}

\subsection{I5: the ledger closes and one inequality is left (T79)}
The transport wall --- the thing that had blocked the line since T20 --- was fully decomposed in T79
(\texttt{transport\_ledger\_probe.py}, $22/22$). The key observation is that the odd-prime side of the
Weil functional agrees \emph{exactly} with the certified plus combination (relative $2.6\times10^{-16}$):
the prime side was never the wall. What remains after the ledger closes is a minimal transport set of
two exact identities (proved), two classically-provable bounds, and exactly one coupled inequality:

\begin{keyeq}$\displaystyle
\textbf{(I5)}\qquad \Qcert(h)+\Delta_2(h)+\Afam(h)-\Ashift(h)\ \ge\ 0
\quad\text{for all autocorrelations }h.$\end{keyeq}

\noindent I5 is the prime\,$\leftrightarrow$\,archimedean coupling. By the closed ledger it is, as an
identity, equivalent to Weil positivity and hence to RH. The diary is explicit about what that means:
\emph{the built-in safeguard fired}. RH does not follow from finite bookkeeping; the hidden infinite
step is exactly I5 --- and the value of T79 is that the step was \emph{found and named} rather than
concealed. This is an equivalence typing, not progress.

T82 then changed I5's \emph{type} (\texttt{heat\_arch\_probe.py}, $22/22$, \textsc{arch-internal}):
$\Darch$ is exactly the internal $\Gamma$ difference $\Afam-\Ashift$ via Legendre duplication, verified
on $18/18$ battery rows at max relative $6.6\times10^{-15}$. I5 stopped being a coupling between two
worlds and became a \emph{self-consistency statement about one heat family}. Again: a type change is
not a proof.

T83 (\texttt{fe\_symmetrization\_probe.py}, $27/27$, \textsc{invariant-null}) then asked the obvious
question of the self-duality picture (Section~\ref{sec:bigpicture}): can the functional equation itself
do the work? It cannot, and the way it fails is informative. The $\zeta$ reflection $J_{1/2}$ acts as
the identity on autocorrelations --- the symmetry is already absorbed in evenness, and
$Q_{\mathrm{Weil}}(Jh)=Q_{\mathrm{Weil}}(h)$ holds floating-point-exactly on $15$ test functions --- and
the family reflection $J_1$ moves nothing in the value geography ($5/24$ overlap invariant, $\lambda^*$
covariant-exact) while pushing $m_+h$ out of the Weil cone in closed form, with
$\mathrm{Fix}(J_1)\cap\{\text{Weil cone}\}=\{0\}$ exactly. The structural finding underneath the null is
the one quoted in the big picture: the product of the two reflections is the unit-line shift, i.e.\ the
transport operator, and the wall is transversal to the functional equation rather than positioned by it.
Symmetrisation is therefore not a route; the same mirror returns, as an obstruction, in T106.

\subsection{The Connes dictionary (T87) and where positivity is anchored}
T87 (\texttt{i5\_connes\_dictionary\_probe.py}, $22/22$, \textsc{dictionary-exact-core}) identified the
I5 one-family form as a compiler instance of the semi-local Connes structure. Four exact matches:
the two-shell orbit route equals Connes' expression at relative $6.3\times10^{-16}$; the orbit kernel of
Connes--Consani (2021) is literally the compiler's Poisson ladder at $\rho=e^u$ (\texttt{sympy} plus
$1.7\times10^{-50}$); $h_+$ is by definition the compiler kernel $k_\zeta=K_{\mathrm{fam}}-K_{\mathrm{shift}}$;
and $k_\zeta=2\theta'_{\mathrm{RS}}$ --- the derivative of the \emph{Riemann--Siegel phase}, at relative
$0.0$. The dihedral shift is the scaling group $\R^*_+$.

The most useful part of T87 is not a match but a \emph{boundary}. The two positivity programs are
complementary: Connes--Consani anchor positivity \emph{before} the primes (the prime-free Sonin window,
support $(\tfrac12,2)$, i.e.\ $u\in(-\log2,\log2)$), the compiler anchors it \emph{after} the primes
(atom certificates, all supports). The window boundary is exactly the first prime atom $u=\log2$
(\texttt{sympy}-exact; on narrow rows the prime side, $\Qcert$ and $\Delta_2$ all vanish identically).
The sectors \emph{touch} and do not overlap; the I5 coupling lives in the crossing region between them,
outside both proved sectors.

\subsection{The tightness map (T88) and the crossing (T89--T91)}
T88 (\texttt{i5\_tightness\_probe.py}, $27/27$) mapped where I5 is actually tight: a band-limited zero
plateau (classical), safe zones, and \emph{eight} near-vertical tight curves whose spacing follows the
$\Gamma$-side density law (ratio $1.01\pm0.08$, computed zero-free --- no spectral identification is
claimed or made). There are no negatives on genuine autocorrelations; the earlier T76 negatives were
the shadow of missing archimedean/$p=2$ bookkeeping, exactly as the ledger predicted.

T89 (\texttt{sonin\_crossing\_probe.py}, $19/19$) established that the window compression of the Weil
form \emph{is} Bombieri's classical object, and settled the sharpness question: the boundary
at $\log2$ is \emph{not} sharp. The margin falls smoothly from $a\approx0.5$ towards the floor with no
structure at $\log2$; the prime-free sector still has margin $0.557$ at its own boundary and keeps
$\ge0.25$ of the reference comfort to $a\approx0.75$ ($+0.057$ beyond the proved boundary). The
residual subspace grows softly $0\to1\to2\to3\to4$ (out of $32$) and carries the global minimum.
T90 (\texttt{residual\_dissection\_probe.py}, $17/17$) made those residual vectors explicit --- $n$-stable
to $\le2.1^\circ$, closed Gaussian$\times\cos$/Gaussian$\times\sin$ fits with $99{+}\%$ capture --- and
decided the core question: the residual is \emph{not} just pole coupling ($3/10$ vectors are controlled
by no structure at all; beyond $a\approx0.85$ genuinely pole-free atom\,$\leftrightarrow$\,arch content
remains).

T91 (\texttt{thin\_band\_analytic\_probe.py}, $19/19$) identified the band as the \emph{one-atom zone}
and found the mechanism that dominates the rest of the document: beyond an inner edge the prime atom
$u=\log2$ becomes \emph{load-bearing} --- the prime \emph{rescues} the form where the archimedean margin
is exhausted (rescue term $+0.22\ldots+0.41$, purely odd sector). Two uncertainty constants were
decided in closed form: $a\cdot t_{\mathrm{rms}}\to\pi$ (Wirtinger, $0.3\%$) and
$a\cdot t_{\mathrm{cent}}\to2\,\mathrm{Si}(\pi)-4/\pi=2.4306$ (two independent routes, $10^{-25}$).

The band that these three parts map out is narrow and its landmarks are explicit. Reading outwards from
the classical side, the map is: classical edge $0.34657$, prime-free crossing $0.3723$--$0.3763$, band
top $\alpha^*=0.46265$, second atom $\log3/2=0.54931$. Everything in Sections~\ref{sec:handover} and
\ref{sec:certification} happens inside that interval, between the first and the second prime atom, and
the four numbers are the coordinates used throughout. The map was re-derived by a third independent
implementation (T93, \texttt{band\_reconciliation\_probe.py}, $41/41$) before it was used further.

An attempt to convert the band into a proof was made and did not succeed. T92
(\texttt{zone\_extension\_proof\_probe.py}, $36/36$, \textsc{t-skeleton}) certified something small and
real --- $Q\ge6.7\times10^{-12}>0$ on an $8$-dimensional sine-window subspace across the whole band,
with an error certificate and a $1057$-point covering --- but not the lattice zone extension it was
aiming at, and it named two independent walls: $\lambda_{\min}$ collapses geometrically (factor $11.4$
per mode), and the high-mode complement would need $N^*\gtrsim5069$ modes. The finding, and the reason
it is quoted here rather than buried, is a statement about instruments: \emph{the finite-block route is
structurally the wrong instrument for this target, not an under-resourced one}. That is what motivated
the change of approach in Section~\ref{sec:handover}.

\subsection{The blind prime demonstration (T94)}
A separate, deliberately modest exercise. T94 (\texttt{prime\_blind\_demo\_probe.py}, $16/16$,
\textsc{blind-100}) used the exact geometric criterion ``odd $n$ is prime $\iff r_4(n)=8(n+1)$''
(Jacobi 1834, on the $\Z^4=\text{rank-}|\mu_4|$ theta tower) to predict all $753$ primes of the
never-seen window $[1\,000\,001,\,1\,010\,000]$ --- count and MD5 declared \emph{before} any ground
truth, then evaluated: exactly $753$, zero false positives, zero false negatives, $100.00\%$ on all
$10^4$ integers, with an AST-enforced prediction path containing no divisibility test (no
\texttt{isprime}, no sieve, no \texttt{\%}, \texttt{//} or \texttt{/}). The cost is part of the
result: a sieve wins by ${\sim}820\times$ ($1$\,ms against $0.83$\,s). This is structure completeness,
not an algorithmic advance --- and the \emph{spectral} prediction (zeros $\to$ primes) remains bound to
I5/RH and is not claimed.

% =====================================================================
\section{The handover induction (T95--T101)}\label{sec:handover}
% =====================================================================
The band work left a picture: each prime atom rescues its own zone and hands over to the next. T95--T101
turned that picture into an induction, and then measured whether the induction can run.

\subsection{The relay is real (T95, T96)}
T95 (\texttt{continuum\_extension\_probe.py}, $28/28$) proved the C1 chain unconditionally
($|h_f(\log2)|\le1/2$ by disjoint supports with margin $0.2305$; the windowed symmetrised shift has
$\|S\|=1/2$ exactly, $1.1\times10^{-16}$) and found that the binding mechanism is \emph{not} the
two-bump extremal (share $0.046$) but the atom rescue: above $\alpha_{\mathrm{free}}=0.371654$ the
prime-free minimiser has $h(\log2)<0$ down to $-0.328$, so the atom adds positivity exactly in the
losing direction.

T96 (\texttt{relay\_test\_probe.py}, $21/21$) is the model of a self-correcting run. The apparent
``edge'' of T95 at $\alpha^*$ was a map artefact: the value $+7.8\times10^{-6}$ reproduced exactly, but
it sits on a smooth, strictly positive curve with no feature at all ($\lambda_{\min}>0$ on the whole of
$[0.38,0.86]$, decay law $\lambda_\infty\sim\exp(-49.2\alpha)$). The \emph{relay}, however, is real ---
and it lives entirely in the counterfactual: remove the $\log3$ atom and $\lambda_{\min}$ crashes to
$-0.445$, with the loser being exactly the anti-double-bump at distance $\log3$ (alignment $-0.99$) and
a rescue identity at $5\times10^{-15}$. All four handover windows are positive: each atom arrives
strictly \emph{before} it is needed ($+0.025/+0.009/+0.011/+0.007$ for $k=1..4$). Three consequences
were recorded: it is a margin problem, not an edge problem; numerics as a witness is exhausted beyond
$\alpha\approx0.55$; and the counterfactual is the proof target, because it works on $O(0.1)$-sized
quantities instead of $10^{-6}$.

\subsection{The induction step and its one missing lemma (T97, T98)}
T97 (\texttt{relay\_induction\_probe.py}, $105/105$) gave the induction step proof shape. Alignment is
true and sharp with a constant law; the $t=0$ killer loss is proved by three structure identities
(old atoms vanish identically on $E_-$; the pole rescue flips with factor $1-\cosh(u/2)$;
$|\hat f|^2=2(1-\cos(tu))|\hat g|^2$), giving $k_{\mathrm{eff}}=(1-\cos(tu))k(t)$ with infimum
$-1.11\ldots-2.59$ instead of $k(0)=-5.37$.\footnote{In closed form (a T92 side result)
$k(0)=-\gamma-3\log2-\pi/2-\log\pi$. Prior-art note (literature pass, 2026-07-27): in substance
this is the classical digamma value $\psi(1/4)=-\pi/2-3\log2-\gamma$, shifted by $-\log\pi$, as it
appears in Suzuki 2023 (J.\ London Math.\ Soc.\ 107) --- a good consistency check, not a find.} The structural pearl: $E_0$ is \emph{literally} the same
form at a smaller window ($7\times10^{-14}$) --- the induction hypothesis, by self-similarity. What
remained was a single cross-block lemma about a single vector.

T98 (\texttt{cross\_lemma\_probe.py}, $44/44$) attacked it and produced the most instructive negative
result of the series: the measured $\sqrt{\lambda_0}$ law is Douglas (1966) range inclusion, i.e.\ a
\emph{tautology} of the positivity one is trying to prove --- the identification \emph{dissolves} the
lemma rather than proving it. Three T97 premises were refuted in the same run. The replacement target
is an exact scalar inequality with no constants and no near-null vector:

\begin{keyeq}$\displaystyle
D_k(\alpha)\ :=\ -\lambda_{\min}\!\left(Q|_{E_-}-Q_{-0}(Q|_{E_0})^{-1}Q_{0-}\right)\ \le\ \tfrac{\mu_k}{2}.$\end{keyeq}

\noindent T99 (\texttt{dk\_recursion\_probe.py}, $23/23$) then found the a-priori structure: an exact
mirror-parity selection rule ($J_-Q_{-0}J_0=-Q_{-0}$ at $2\times10^{-15}$) giving
$D_k=\max(D_{\mathrm{even}},D_{\mathrm{odd}})$, with the binding even channel coupling only to the
\emph{odd} spectrum of the smaller window --- so the fragile near-null mode is identically absent from
the binding channel. T100 (\texttt{bulk\_remainder\_probe.py}, $27/27$) closed the bulk remainder from
$11/24$ to $24/24$ samples and showed the $10.8\%$ drift was a grid artefact.

\subsection{The race (T101)}
T101 (\texttt{k\_asymptotics\_probe.py}, $31/31$, \textsc{crowding-trends}) resolved $16$ zones
($n=2\ldots29$) and separated two questions that had been fused. The \emph{mathematics} trends
self-supportingly: the primitive quantities are flat over $16$ zones ($w/g$ slope $+0.002\pm0.178$;
the relative $D_k$ margin does not trend; $D_k\le\mu_k/2$ holds $64/64$ and never fails). Atom
\emph{strengths} do not fall ($\mu\sim n^{-0.087}$), only the \emph{spacings} do ($n^{-0.615}$, exact) ---
pure arithmetic. What loses the race is the \emph{closure instrument}: the demand ratio $r_k$ decays
like $\exp(-0.1622k)$ (fit, $\pm0.0562$). The collapsed one-parameter law
$w_k=0.0838\cdot(\text{atom spacing})/\mu_k$ reduced the scatter from $8.85\times$ to $2.66\times$ with
zero residual trend.

\begin{figure}[ht]\centering
\begin{tikzpicture}[font=\small,x=0.52cm,y=2.6cm]
% axes
\draw[->,tfgray!80,thick] (0,-0.25) -- (17.4,-0.25) node[right,black]{zone $k$};
\draw[->,tfgray!80,thick] (0,-0.25) -- (0,1.15) node[above,black]{$\log_{10} r_k$};
\foreach \k in {1,4,8,12,16} \draw[tfgray!60] (\k,-0.27) -- (\k,-0.23) node[below=2pt,black,font=\scriptsize]{\k};
\foreach \y in {0,0.5,1.0} \draw[tfgray!60] (-0.25,\y) -- (0.25,\y) node[left=6pt,black,font=\scriptsize]{$\y$};
% spectrum floor (schematic)
\fill[tfred!7] (0,-0.25) rectangle (17.4,0.0);
\draw[tfred!60!black,dashed] (0,0.0) -- (17.4,0.0);
% T101 instrument: slope -0.1622 per zone in ln  ->  -0.07044 in log10
\draw[tfred!75!black,very thick] (1,0.970) -- (16,-0.087);
% T103 instrument: 9.33 -> 2.70
\draw[tfgreen!60!black,very thick] (1,0.970) -- (16,0.431);
\filldraw[tfgreen!60!black] (1,0.970) circle (1.4pt);
\filldraw[tfgreen!60!black] (16,0.431) circle (1.4pt);
\filldraw[tfred!75!black] (16,-0.087) circle (1.4pt);
% labels last, on an opaque backing so no rule ever crosses a glyph
\node[tfred!60!black,font=\scriptsize,anchor=west,inner sep=1.5pt]
  at (0.4,-0.13) {spectrum floor (schematic)};
\node[tfred!75!black,font=\scriptsize,align=left,anchor=west,fill=white,inner sep=2pt]
  at (2.6,0.25) {T101 instrument\\[-1pt]slope $-0.1622\pm0.0562$};
\node[tfgreen!45!black,font=\scriptsize,align=left,anchor=west,fill=white,inner sep=2pt]
  at (9.2,0.90) {T103 instrument\\[-1pt]slope $-0.0748\pm0.0116$};
\node[font=\scriptsize,anchor=south west,inner sep=1.5pt] at (1.15,0.99) {$9.33$};
\node[font=\scriptsize,anchor=west,inner sep=2.5pt] at (16.1,0.431) {$2.70$};
\end{tikzpicture}
\caption{\schem The instrument race over the $16$ resolved zones. \emph{Literal:} the two decay slopes
($-0.1622\pm0.0562$ measured in T101; $-0.0748\pm0.0116$ measured in T103, a factor $2.2$ flatter), and
the T103 endpoints $r_1=9.33\to r_{16}=2.70$ with the statement that the improved instrument never
leaves the spectrum. \emph{Schematic:} the straight-line interpolation in $\log_{10}$, the common
anchor of the two lines, and the position of the ``spectrum floor''. Both slopes are declared fits in
the diary. What loses the race is the instrument, not the inequality: the true inequality
$D_k\le\mu_k/2$ holds $64/64$ and its relative margin does not trend.}\label{fig:race}
\end{figure}

\noindent T101 also stated what an asymptotic theorem would need, and that list has not changed since:
(A) the collapsed lower bound $w_k\ge c\,g_k/\mu_k$ with a fixed $c>0$ --- a sharp, non-perturbative,
\emph{arithmetic} lower bound at a Weil-positivity edge, with Taylor arguments forbidden by the
$\exp(-c/\delta')$ onset; (B) a uniform relative margin; (C) a replacement for the bulk remainder;
(D) a finite check of the low zones. If crowding ever wins, the hardness sits entirely in (A).

% =====================================================================
\section{The certification arc (T102--T115)}\label{sec:certification}
% =====================================================================
The last fourteen parts stopped improving measurements and started converting them into certificates ---
and, where that failed, into explicit refutations.

\subsection{The mechanism: a crossing, not a singularity (T102)}
T102 (\texttt{arithmetic\_bound\_probe.py}, $42/42$, \textsc{mechanism-identified}) computed the
mechanism exactly. The $k$-th atom acts on $E_-/E_0/E_+$ as $\mathrm{diag}(-\tfrac12,0,+\tfrac12)$, so
$Q_{k-1}=\Qfull-\tfrac{\mu_k}{2}P_-+\tfrac{\mu_k}{2}P_+$: the atom strength $\mu_k$ enters exactly once,
linearly. That gives a two-sided sandwich through the Schur profile $\sigma_k(\delta)$, with
$\mu_k/2\le\sigma_k$ implying that the handover holds, and $2w_k$ lying between the crossings in
$16/16$ zones.

Two corrections came out of the same run. First, the onset is \emph{anchored} at $\delta_c=2w_k$ --- a
crossing of two finite quantities ($R^2=0.968$) --- and exponential-singular versus power behaviour is
not separable; T96's essential-singularity reading is compatible but no longer distinguished. Second,
and more consequentially, the binding constraint \emph{flips}: no concentration condition binds. The
bare $E_-$ form is strongly positive ($2.65\ldots3.52$, i.e.\ $4$--$14\times$ $\mu_k/2$), the pole flip
and band mass saturate their classical ceilings ($96.9$--$97.1\%$ of Cauchy--Schwarz and
Landau--Pollak/prolate respectively), and the onset is manufactured \emph{entirely} by the Schur
dressing against $E_0\oplus E_+$, which eats $35.7\%$--$97.3\%$ of the bare eigenvalue. The loser does
not live in $E_-$.

\begin{figure}[ht]\centering
\begin{tikzpicture}[font=\small,x=1.0cm,y=1.0cm]
\draw[->,tfgray!80,thick] (0,0) -- (10.2,0) node[right,black]{$\delta$};
\draw[->,tfgray!80,thick] (0,0) -- (0,4.1) node[above,black]{};
% flat mu_k/2 line
\draw[tfblue!70!black,very thick] (0,1.55) -- (10,1.55);
\node[tfblue!70!black,anchor=west,font=\scriptsize] at (8.1,1.85) {$\mu_k/2$ (flat)};
% upper and lower schur profile brackets
\draw[tfgreen!55!black,very thick] plot[smooth,tension=0.7] coordinates {(0.4,3.75) (2.2,3.1) (4.0,2.3) (5.9,1.55) (7.6,1.02) (9.4,0.72)};
\draw[tfgreen!55!black,very thick,dash pattern=on 3pt off 2pt] plot[smooth,tension=0.7] coordinates {(0.4,3.15) (2.0,2.5) (3.4,1.85) (4.3,1.55) (6.0,1.02) (8.4,0.66)};
\node[tfgreen!45!black,anchor=west,font=\scriptsize,align=left,fill=white,inner sep=2pt] at (5.9,3.72)
  {Schur profile $\sigma_k(\delta)$\\[-1pt](two-sided bracket)};
% crossings
\filldraw[tfgreen!45!black] (4.3,1.55) circle (1.6pt);
\filldraw[tfgreen!45!black] (5.9,1.55) circle (1.6pt);
\draw[tfgray!70,dashed] (4.3,0) -- (4.3,1.55);
\draw[tfgray!70,dashed] (5.9,0) -- (5.9,1.55);
\draw[decorate,decoration={brace,amplitude=4pt,mirror},tfgold!70!black,thick]
  (4.3,-0.12) -- (5.9,-0.12) node[midway,below=4pt,font=\scriptsize,black]{$2w_k$ (handover window)};
\node[anchor=south,font=\scriptsize,tfgray] at (2.1,0.15) {handover holds: $\sigma_k\ge\mu_k/2$};
\node[anchor=south,font=\scriptsize,tfgray] at (8.3,0.15) {atom $k{+}1$ takes over};
\node[anchor=west,font=\scriptsize,tfred!70!black,align=left,fill=white,inner sep=2pt] at (6.15,2.45)
  {onset anchored at $\delta_c=2w_k$\\[-1pt]($R^2=0.968$) --- a crossing,\\[-1pt]not a singularity};
\end{tikzpicture}
\caption{\schem The crossing mechanism of T102. \emph{Literal:} $\mu_k$ enters exactly once and
linearly, so the condition is a two-sided sandwich through the Schur profile $\sigma_k(\delta)$ against
the flat level $\mu_k/2$; the window width $2w_k$ lies between the crossings in $16/16$ zones; the onset
is anchored at $\delta_c=2w_k$ with $R^2=0.968$; the candidate bound $w_k\ge A_k\mu_k^{1/q_k}$ holds
$16/16$ with ratio $0.749$--$0.940$. \emph{Schematic:} the curve shapes and all absolute
scales.}\label{fig:crossing}
\end{figure}

\noindent The same run performed the most useful kill of the arc: a triple decomposition test showed
that the T101 law in $C/g$ was a \emph{proxy}. $g_k$ is causally impossible (the form $Q_{k-1}$ is blind
to the next atom), statistically dispensable (a causal law $C\cdot u$ has $\mathrm{CV}$ $23.3\%$ against
$27.6\%$), and arithmetically only a ceiling ($C_k\le g_k\mu_k$ exactly; the $16$-zone extrapolation
violates the ceiling from $k=69$; checked over $18\,120$ prime-power atoms to $n=200\,000$).

\subsection{The instrument, rebuilt (T103), and the two doors}
T103 (\texttt{instrument\_probe.py}, $29/29$, \textsc{instrument-improved}) rebuilt the closure
instrument with certified weights $w_j\ge1/\lambda_j$ and a $\theta$-weighted band sum, reproducing both
anchors exactly. With \emph{one} fixed, $k$-uniform instrument ($\Lambda_0=3$, $r=2$, no per-zone
tuning) the closure map jumps from $7/64$ to $44/64$ and the race slope halves
(Figure~\ref{fig:race}). The price is explicit: the number of explicit modes grows $2\to232$.
And the verdict relocated the remainder: it is no longer the bulk (the bulk is measurably not
low-rank --- effective rank up to $0.579\dim E_-$) but the \emph{wing slack} $S=1-\rho$, which falls
$0.2091\to0.0392$.

\begin{figure}[ht]\centering
\begin{tikzpicture}[font=\scriptsize,x=1cm,y=1cm]
\def\cw{0.30}
% left grid: 7/64
\begin{scope}[shift={(0,0)}]
  \foreach \c in {1,...,16}{\foreach \r in {1,...,4}{
     \draw[tfgray!35] ({(\c-1)*\cw},{(\r-1)*\cw}) rectangle ++(\cw,\cw);}}
  \foreach \c/\rmax in {1/4,2/2,3/1}{\foreach \r in {1,...,\rmax}{
     \fill[tfgreen!45!black,opacity=0.75] ({(\c-1)*\cw+0.02},{(\r-1)*\cw+0.02}) rectangle ++({\cw-0.04},{\cw-0.04});}}
  \node[anchor=north west,align=left,font=\scriptsize] at (0,-0.12)
    {\textbf{T101 instrument: 7/64}\\[-1pt]closes only in the first zones};
  \node[rotate=90,anchor=south,font=\scriptsize,tfgray] at (-0.28,0.6) {wing fraction};
\end{scope}
% right grid: 44/64
\begin{scope}[shift={(6.1,0)}]
  \foreach \c in {1,...,16}{\foreach \r in {1,...,4}{
     \draw[tfgray!35] ({(\c-1)*\cw},{(\r-1)*\cw}) rectangle ++(\cw,\cw);}}
  \foreach \c/\rmax in {1/4,2/4,3/4,4/4,5/4,6/4,7/4,8/4,9/4,10/3,11/3,12/2}{\foreach \r in {1,...,\rmax}{
     \fill[tfgreen!45!black,opacity=0.75] ({(\c-1)*\cw+0.02},{(\r-1)*\cw+0.02}) rectangle ++({\cw-0.04},{\cw-0.04});}}
  \node[anchor=north west,align=left,font=\scriptsize] at (0,-0.12)
    {\textbf{T103 instrument: 44/64}\\[-1pt]one fixed $k$-uniform instrument ($\Lambda_0=3$, $r=2$)};
\end{scope}
\node[anchor=south,font=\scriptsize,tfgray] at (2.4,1.28) {$16$ zones $\times$ $4$ wing fractions};
\node[anchor=south,font=\scriptsize,tfgray] at (8.5,1.28) {$16$ zones $\times$ $4$ wing fractions};
\end{tikzpicture}
\caption{\schem The closure map. \emph{Literal:} the grid is $16$ zones $\times$ $4$ wing fractions
$=64$ cells; the old instrument closes $7$ of them, the rebuilt one closes $44$ with a single fixed
$k$-uniform setting ($\Lambda_0=3$, $r=2$) and no per-zone tuning. \emph{Schematic:} \emph{which} cells
are shaded --- the diary records the totals, not the pattern; the shading is drawn as a monotone wedge
purely to visualise the counts.}\label{fig:closuremap}
\end{figure}

\noindent The convergence recorded at the end of T102 is the pivot of the whole arc: T103's wing slack
$S=1-\rho$ and T102's Schur dressing are \emph{the same object seen from two sides}. Both doors lead to
the wing near-null direction.

\begin{figure}[ht]\centering
\begin{tikzpicture}[font=\small,
  box/.style={draw=tfblue!70!black,fill=tfblue!5,rounded corners=2pt,inner sep=5pt,align=center,
              text width=64mm,minimum height=18mm},
  tgt/.style={box,draw=tfgold!70!black,fill=tfgold!8,text width=86mm,minimum height=15mm},
  nxt/.style={box,draw=tfgreen!65!black,fill=tfgreen!6,text width=86mm,minimum height=13mm},
  ar/.style={-{Latex[length=2.2mm]},draw=tfgray!80,thick}]
\node[box] (a) at (-3.9,0) {\textbf{Door A} --- T102 (arithmetic bound)\\[2pt]
  \footnotesize the onset is manufactured by the Schur dressing against $E_0\oplus E_+$;
  it eats $35.7\%$--$97.3\%$ of the bare eigenvalue};
\node[box] (c) at (3.9,0) {\textbf{Door C} --- T103 (instrument)\\[2pt]
  \footnotesize the remainder is not the bulk but the wing slack $S=1-\rho$,
  falling $0.2091\to0.0392$};
\node[tgt] (o) at (0,-3.4) {\textbf{One object:} the wing near-null direction\\[2pt]
  \footnotesize dressing almost rank one (effective rank $1.07$--$2.7$; $99.7$--$100\%$ of the
  near-null direction in \emph{one} prolate mode)};
\node[nxt] (n) at (0,-6.1) {T104 $\to$ T105 $\to$ T106 $\to$ T107 $\to$ T108 $\to$ T109 $\to$ T110 $\to$ T111 $\to$ T112 $\to$ T113 $\to$ T114 $\to$ T115\\[2pt]
  \footnotesize certify \emph{bare}, prove the avoidance law, split the parity channels,
  reduce to one scalar ratio, then to two scalars, certify both --- then propagate the margin,
  measure where the propagation fails, rescale the frame, settle the currency question,
  dissolve the wall --- then split the transport and break the compute cap};
\draw[ar] (a) -- (o);
\draw[ar] (c) -- (o);
\draw[ar] (o) -- (n);
\end{tikzpicture}
\caption{Two doors, one object. Both entries are literal diary findings; the convergence statement
(``the same object from two sides'') is quoted from the T102 closing note.}\label{fig:doors}
\end{figure}

\subsection{Inductive chains, and the death of the margin route (T104)}
T104 ran as \emph{two independent arms} after a file collision between two parallel workers was resolved
by keeping both implementations of the same contract --- \texttt{schur\_profile\_bound\_probe.py}
($21/21$) and \texttt{schur\_profile\_chain\_probe.py} ($47/47$). All core findings converged, which
makes it a de-facto independent replication.

\begin{itemize}[leftmargin=1.5em]
\item \textbf{The naive margin route is dead}: $M\ge m\cdot\mathrm{Id}$ holds in $0/16$ zones; the
coupling mass to margin ratio is $10^3$--$10^6$; $m$ vanishes in the continuum like $M^{-1.7}$ (fit).
\item \textbf{The mechanism is avoidance}: the coupling $B_-$ decouples from the soft bulk modes ---
the softest mode carries $0.00\%$ of the coupling mass for $n\ge3$, and the mass sits in mid/high
bands (lowest band only $2$--$12\%$). Positivity is not saved by a margin but by structural avoidance.
\item \textbf{Exact spectral-split chains close $16/16$} with finite data: arm A via the exact
block-inverse identity with matrix caps in PSD order (headroom $1.15$--$1.44$, reach
$\gamma_{\max}=0.50$--$0.90$); arm B via a threshold split with an $O(1)$, resolution-stable threshold
$w=2.47$--$5.90$ --- but requiring $r_{\min}=64\ldots1024$ soft modes $\propto M$, which makes the
continuum statement a spectral-density condition.
\item A recommendation from T96/T103 was tested and \emph{rejected}: the prolate wing basis loses
against the raw basis in both arms.
\end{itemize}

\noindent Consequence: the hard core migrated from $\sigma_k$ as a whole to four named ingredients ---
a lower bound on $\mathrm{bare}_k=\lambda_{\min}(\Qfull|_{E_-})$, the soft dressing scalar $L$, finite
induction data, and a one-sided edge estimate instead of a fit. The chain around them is classical
(Schur complement, exact block inverse, Parseval/Bessel, Weyl) and fit-free.

\subsection{The certified bare bound and the avoidance theorem (T105)}
T105 (\texttt{bare\_wing\_bound\_probe.py}, $28/28$, \textsc{one-of-two}) delivered one of the two
missing ingredients in closed form. First it resolved the T104 arm discrepancy: \emph{bare} depends only
on $(u_k,\delta)$ --- the window centre cancels ($\le0.74\%$ drift under $8\times$ refinement) --- with
the exact split $\mathrm{bare}=\mu_k/2+b_0$. Then three classical steps (Bessel; a Legendre/level-set
optimisation at $k(\pi/\delta)$; Cauchy--Schwarz at the pole pair) collapse to a closed lower bound:

\begin{keyeq}$\displaystyle
b_0\ \ge\ \frac{\delta}{\pi}\int_0^{\pi/\delta}\!\!k(t)\,dt\;-\;\delta\,K_{\max}\;-\;4\left(\cosh\tfrac u2-1\right)\sinh\tfrac\delta2 .$\end{keyeq}

\noindent This is the band average of the archimedean kernel over the band conjugate to the wing, of
size ${\sim}\log(1/2\delta)-1$. It is positive in $16/16$ zones ($1.27$--$3.54$) at $81.1$--$92.7\%$ of
the measured value (median $92\%$), it sharpens as the wing flattens, and it needs no eigenvalue or
induction data.

Second, the \emph{avoidance law} became a theorem rather than an observation:
$\Qfull\succeq0\Rightarrow\|B^{\mathsf T}v_i\|^2\le w_i\Lambda_-$ (a semi-inner-product argument;
coupling proportional to the mode eigenvalue), verified on every mode of every zone with worst ratio
$0.377$. That explains the measured $0.00\%$ avoidance exactly. Alongside it, an exact parity
superselection: $JQJ=Q$ ($5.7\times10^{-15}$) and $JB=-BR$ ($3.4\times10^{-15}$) --- the two channels do
not mix, and the softest mode is $J$-even in $16/16$ zones. Notably the \emph{Euclidean} angles are not
small (cosine up to $1.0$): the avoidance lives in the form metric, as a Friedrichs angle.

What T105 left is one Loewner statement:
$\Qfull(\alpha)\succeq(\mu_k/2)P_-^{(k)}\iff\rho\le1-(\mu_k/2)/\mathrm{bare}_k$ --- a Friedrichs-angle /
spectral-density condition, explicitly \emph{not} a margin condition. The free cap
$L\le\Lambda_-=5.63$--$7.52$ diverges like $\log(1/D)$, so the certificate has to come from the spectral
density near the soft end.

\subsection{Parity superselection: the pole splits (T106)}
T106 (\texttt{friedrichs\_angle\_probe.py}, $32/32$, \textsc{density-mapped}) killed two routes and
found the structural result of the arc. The kills are in Section~\ref{sec:kills}. The finding:

\begin{keyeq}$\displaystyle
ab^{\mathsf T}+ba^{\mathsf T}\;=\;ss^{\mathsf T}-tt^{\mathsf T}\qquad(10^{-18};\ Ja=b)$\end{keyeq}

\noindent Here $J$ is the mirror parity of Section~\ref{sec:bigpicture}, arriving for the second time
and now as the obstruction itself. The rank-2 Weil pole splits \emph{exactly} into a positive rank-one lift in the $J$-even
channel and a negative rank-one pressure in the $J$-odd channel. The Toeplitz part has exactly one
negative eigendirection; it is $J$-even, and $ss^{\mathsf T}$ removes it precisely there. The odd
channel never needed the pole at all.

\begin{figure}[ht]\centering
\begin{tikzpicture}[font=\small,
  src/.style={draw=tfgold!70!black,fill=tfgold!8,rounded corners=2pt,inner sep=5pt,align=center,text width=52mm},
  ev/.style={draw=tfgreen!65!black,fill=tfgreen!6,rounded corners=2pt,inner sep=5pt,align=center,text width=62mm},
  od/.style={draw=tfred!70!black,fill=tfred!5,rounded corners=2pt,inner sep=5pt,align=center,text width=62mm},
  ar/.style={-{Latex[length=2mm]},draw=tfgray!80,thick}]
\node[src] (p) at (0,0) {rank-2 Weil pole\\[1pt]
  $ab^{\mathsf T}+ba^{\mathsf T}=ss^{\mathsf T}-tt^{\mathsf T}$ \ \footnotesize($10^{-18}$)};
\node[ev] (e) at (-4.3,-2.3) {\textbf{$J$-even channel}\\[1pt]
  \footnotesize $+ss^{\mathsf T}$: rank-one \emph{lift}. Carries the single negative Toeplitz
  direction --- removed exactly there. Soft edge lives here ($3.7$--$100\times$ softer);
  $\rho_+=0.076$--$0.464$. Angle chain closes $16/16$ ($0.107$--$0.583\times$ budget).};
\node[od] (o) at (4.3,-2.3) {\textbf{$J$-odd channel}\\[1pt]
  \footnotesize $-tt^{\mathsf T}$: rank-one \emph{pressure}. Never needed the pole; Toeplitz part
  \emph{positive}. $\rho_-=0.557$--$0.826$ --- the dangerous channel despite the better density.
  Angle chain closes $9/16$.};
\node[draw=tfgray!60,fill=tfgray!4,rounded corners=2pt,inner sep=5pt,align=center,text width=86mm]
  (r) at (0,-4.85) {\textbf{The remaining object (R)}\\[1pt]
  $\Qfull(\alpha_k)|_{J=-1}\ \succeq\ (\mu_k/2)\,P_-|_{J=-1}$ \quad on $\lceil p/2\rceil$ dimensions};
\draw[ar] (p) -- (e);
\draw[ar] (p) -- (o);
\draw[ar] (o) -- (r);
\end{tikzpicture}
\caption{Parity superselection (T106). All numbers are literal; the layout is illustrative. The
practical consequence: it is not the density that decides the angle but the \emph{avoidance} --- the
dangerous channel is the one with the better density. T105's certified bare bound sits $16/16$ exactly
in the hard channel, and $\rho_-$ is resolution-stable ($\le17.6\%$).}\label{fig:parity}
\end{figure}

\noindent The channel-wise angle chain closes $16/16$ in the even channel and $9/16$ in the odd
channel, against only $3/16$ unsplit; all $16$ handovers close at operating depth in both channels under
the exact Schur criterion, with certified $b_0$ at $89$--$93\%$. Certified after T106: bare, the parity
split, the pole splitting, the channel inertia, the avoidance law, the growth inequality, the
accumulation bounds. Measured: $\rho_-$.

\subsection{One scalar ratio (T107)}
T107 (\texttt{odd\_channel\_closure\_probe.py}, $30/30$, \textsc{scalar-tractable}) reduced the matrix
statement (R) to a scalar one. Sherman--Morrison gives exactly
$\text{(R)}\iff s^*\le1$ (identity at $3.4\times10^{-11}$; equivalence confirmed $16/16$ against
Rayleigh--Ritz; precondition $G\succ0$ via Cholesky $16/16$) --- but $s^*$ is \emph{saturated}
($0.99885\ldots0.99999$, with only $n=2$ at $0.2569$), so no uniform $s^*$ chain can work. The Woodbury
decomposition rescues it: $s^*=\tau+\kappa$ and

\begin{keyeq}$\displaystyle
\text{(R)}\iff r:=\frac{\kappa}{\varepsilon}\le1,\qquad \varepsilon:=1-\tau,
\qquad\text{measured } r=0.0053\ldots0.1814 .$\end{keyeq}

\noindent Two orders of magnitude of headroom instead of five decimal places, and $s^*$ is
resolution-stable (spread $\le2.60\%$ over $M=600/900/1200$). The synthesis is stated with the usual
asymmetry: certified-closed $0/16$, measured $16/16$ (with $q=1.31$--$4.16$ at $M=1200$, replacing
T106's $0.747$--$1.156$ --- a factor $1.7$--$2.5$ gain). Exactly one object is missing: a certified
positive lower bound for $\varepsilon=1-\tilde t^{\mathsf T}T_{\mathrm{odd}}^{-1}\tilde t$ --- how far
the pole vector is from exhausting the odd channel --- or, since $\varepsilon$ and $\kappa$ both fall
with $M$ ($\varepsilon\sim M^{-1.99\ldots-1.69}$, fit), a certificate for the \emph{ratio} $r$ itself.
The named route is via a wing statement of the T105 type plus an avoidance norm, explicitly not via
Szeg\H{o}.

\subsection{The \texorpdfstring{$\varepsilon$}{epsilon} identity, and two scalars (T108)}
T108 (\texttt{ratio\_certificate\_probe.py}, $44/44$, \textsc{epsilon-identity}) turned the missing
positivity of $\varepsilon$ from a bound to be won into an \emph{identity}. Everything hangs on the pole
response $x:=T_{\mathrm{odd}}^{-1}\tilde t$, and three exact identities (verified at
$10^{-11}$--$10^{-13}$) pin it down:
$\varepsilon=x^{\mathsf T}Q|_{\mathrm{odd}}x/\tau$ --- so $\varepsilon$ \emph{is} a
$Q|_{\mathrm{odd}}$-energy, not an accident of subtraction ---
$1/\varepsilon=1+\tilde t^{\mathsf T}Q|_{\mathrm{odd}}^{-1}\tilde t$, and the overlap of $\varepsilon$
and $\kappa$ is exactly one rank-one term $hh^{\mathsf T}/\varepsilon$, so that
$\text{(R)}\iff\hat r=(\mu/2)\lambda_{\max}(\Gamma+hh^{\mathsf T}/\varepsilon)\le1$. The third route
supplies the classical reading: $\tilde t$ is exactly two-term geometry, $\tau$ a $2\times2$
Christoffel--Darboux form, and $\varepsilon$ exactly the square of the last Cholesky pivot --- the
Szeg\H{o}/Levinson prediction error. \emph{The positivity of $\varepsilon$ coincides with the induction
positivity in the odd channel}; it is not a second, independent thing to certify.

Two consequences, and one correction of the previous picture. The consequence: (R) reduces from an
$M/2$-dimensional Loewner statement to \textbf{two scalars} --- $\omega<1$ (measured
$0.2366$--$0.7531$, $16/16$; the T105 structure survives the parity split verbatim,
$V^{\mathsf T}S|_{\mathrm{odd}}V=-\tfrac12 I$ at $4.9\times10^{-14}$) and one statement at the vector
$x$. The chain $\text{(R)}\Leftarrow[\omega<1]\wedge[\lambda_{\min}(Q|_{\mathrm{odd}})\ge
(\mu/2)\tau\theta/(1-\omega)]$ has $C1$ closing $16/16$ (margin $5.5$--$178$; $\ge1$ at all $48$ ladder
points, $2.7$--$343$), $C2$ $15/16$, $C3$ $0/16$ --- the first time a margin chain in this arc is
\emph{not vacuous}: the induction margin it needs is explicitly $9.3\times10^{-8}$--$3.7\times10^{-3}$,
i.e.\ two to six orders of magnitude \emph{below} $\mu/2$. The correction: the soft-edge picture of T107
was wrong in its mechanism --- $\tau$ sits to $99.66\%$ on modes $\lambda\ge1$, and $\|h\|^2$ is small by
sign cancellation ($\sum|k|=401$ versus $\sum k=1$), not by softness. T107's Cauchy--Schwarz chain turns
out to be exactly the Weyl bound (slack $1.0000$).

What is exactly open after T108 is one boundary value. The remaining object is a certified upper bound
for the avoidance norm $\|V^{\mathsf T}T_{\mathrm{odd}}^{-1}\tilde t\|^2$, and on $8$ zones that is
literally the single edge entry $x_0$ of an explicit vector (suppressed at
$|x_0|/\max|x|=2.9\times10^{-3}$--$5.7\times10^{-3}$, with a linear edge profile $r^{1.01}$ quoted as a
fit) --- a boundary-decay statement about an explicit vector, the same object class as the T105 carrier
separation. $\hat r$ is flat in $M$ ($6.1\times$ more stable than $\varepsilon$) and $r,\hat r<1$ at all
$92$ (zone, $M$) points up to $M=3000$, tightest at $0.9648$: measured, on resolved zones, as always.
That boundary value, and the still-uncertified $\omega$, are exactly what the next part attacked.

\subsection{Boundary certificates: cancellation, not decay (T109)}
T109 (\texttt{boundary\_decay\_probe.py}, $29/29$, \textsc{boundary-certified}) closed the arc: both
remaining scalars are now certified, and the whole chain for (R) closes $16/16$ conditional on exactly
one explicit input. Three findings and one honest remainder.

First, the mechanism. The edge value $x_0$ is not small because anything \emph{decays} --- it is small
because a sum \emph{cancels}. Carrier separation is excluded geometrically: $\tilde t$ has its maximum
at $r=0$; the source sits on the boundary itself. The Combes--Thomas hypothesis is satisfied
($T_{\mathrm{odd}}$ is of Jaffard class, with local rate $\rho(s)=D(\tfrac12+2/(e^{2s}-1))$, matched at
a factor $0.85$--$1.00$, quoted as a fit) --- but its conclusion is empty: the Green sum for $x_0$
cancels by a factor $44$--$4.6\times10^{4}$ on $15/16$ zones (the known outlier $n=2$ does not cancel),
and the edge profile is a nearly linear zero, $r^{0.77\ldots1.28}$ as a fit. A boundary condition, not
a decay length --- and no absolute-value bound can survive a cancellation.

Second, the $x_0$ certificate that works. Four candidates were run certified-against-measured:
Combes--Thomas, evaluated with the \emph{exact} Green row so that no constant could rescue it, is
\emph{refuted} ($1/16$; an upper bound on its whole class); a $T$-metric Cauchy--Schwarz at the edge
entry closes $1/16$; the exact Szeg\H{o}/Levinson two-point evaluation exhibits the cancellation but
bounds nothing. The fourth is the new object: a residual / goal-oriented certificate (Prager--Synge;
second order in the two residuals) which \emph{carries} the cancellation by construction instead of
bounding it away. Sharpness $1.0000$--$2.4937$ against the measured norm, $16/16$ within budget.

Third, $\omega$ is cracked --- unconditionally. The certificate is a graded matrix cap in PSD order:
one trial level per soft direction, validity by a single Cholesky factorisation (Sylvester inertia)
plus Loewner antitonicity of the inverse, $n_{\mathrm{top}}=17$--$512$. It gives
$\omega_{\mathrm{cert}}=0.2561$--$0.7569$ against the measured $0.2366$--$0.7531$ --- a factor
$1.004$--$1.124$ --- with \emph{no} hypothesis input. The compression Schur distance that blocked T108
shrinks from $0.18$--$0.52$ to $0.0024$--$0.039$, and the ungraded cap (one level for the whole soft
block, which is what T108's global route amounted to) is vacuous everywhere: the grading is the whole
difference. The trial data are numerically generated; their \emph{validity} is trial-independent
(identity $+$ Cholesky $+$ Loewner), only their sharpness is not.

The synthesis feeds both certificates from one object
($\Gamma_{\mathrm{cert}}=S_{\mathrm{cert}}^{-1}$): the chain closes $16/16$ with margin $1.3$--$10.4$
and is $M$-stable ($44/48$ ladder points to $M=3000$; the certification price $1.5$--$28.6$ does not
grow). What remains open is exactly \textbf{one input}: the \emph{strict} positivity margin of the odd
channel, $\lambda_{\min}(Q|_{\mathrm{odd}})\ge1.7\times10^{-6}$--$7.5\times10^{-3}$ against a measured
$1.4\times10^{-5}$--$5.1\times10^{-2}$ --- two to six orders of magnitude \emph{below} $\mu/2$, i.e.\
strictly weaker than the conclusion it buys; the induction hypothesis as it stands supplies only the
non-strict $Q|_{\mathrm{odd}}\succeq0$. Declared caveats: floating point is not audited, and
$\lambda_{\min}$ is read through Rayleigh--Ritz. Whether that one number propagates itself through
the induction is exactly what the next part measured.

\begin{figure}[htp]\centering
\begin{tikzpicture}[font=\small,
  eq/.style={draw=tfgold!70!black,fill=tfgold!8,rounded corners=2pt,inner sep=4.5pt,align=center,text width=104mm},
  st/.style={draw=tfblue!70!black,fill=tfblue!5,rounded corners=2pt,inner sep=4.5pt,align=center,text width=104mm},
  op/.style={draw=tfred!70!black,fill=tfred!5,rounded corners=2pt,inner sep=4.5pt,align=center,text width=104mm},
  ar/.style={-{Latex[length=2.2mm]},draw=tfgray!85,thick},
  lbl/.style={font=\scriptsize,tfgray,align=left,anchor=west}]
\node[eq] (rh) at (0,0) {\textbf{RH} $\iff$ Weil positivity $\iff$ \textbf{I5}
  \ \footnotesize(one-family form; T79/T82, ledger \vref{v541})};
\node[st] (hand) at (0,-1.55) {per-zone \textbf{handover induction}: \ $D_k(\alpha)\le\mu_k/2$
  \ \footnotesize(T98; holds $64/64$ on the $16$ resolved zones)};
\node[st] (schur) at (0,-3.1) {\textbf{Schur profile} sandwich: \ $\sigma_k(\delta)\ge\mu_k/2$
  \ \footnotesize(T102; $\mu_k$ linear, onset anchored at $\delta_c=2w_k$)};
\node[st] (bare) at (0,-4.65) {$\mathrm{bare}_k$ \textbf{certified} $+$ soft dressing $L$
  \ \footnotesize(T105; $b_0$ closed form, positive $16/16$ at $81.1$--$92.7\%$)};
\node[st] (par) at (0,-6.2) {\textbf{parity split}: even channel closes $16/16$; hardness $\to$ odd
  channel \footnotesize(T106)};
\node[st] (loew) at (0,-7.75) {\textbf{(R)}: $\Qfull|_{J=-1}\succeq(\mu_k/2)P_-|_{J=-1}$ on
  $\lceil p/2\rceil$ dimensions \footnotesize(T106)};
\node[st] (ratio) at (0,-9.3) {$\iff s^*\le1$ (saturated) $\iff$ \ \textbf{$r=\kappa/\varepsilon\le1$}
  \ \footnotesize(T107; measured $0.0053$--$0.1814$)};
\node[st] (eps) at (0,-10.85) {$\iff$ two scalars: $\omega<1$ \ and one statement at
  $x=T_{\mathrm{odd}}^{-1}\tilde t$ \ \footnotesize(T108; $\varepsilon$ identity; boundary value $x_0$)};
\node[st] (bnd) at (0,-12.4) {\textbf{both scalars certified}: $\omega_{\mathrm{cert}}<1$
  unconditional (graded matrix cap) $+$ residual certificate carrying the boundary cancellation
  \ \footnotesize(T109; closes $16/16$, margin $1.3$--$10.4$)};
\node[op] (marg) at (0,-14.05) {open \textbf{input}: one strict-positivity margin
  $\lambda_{\min}(Q|_{\mathrm{odd}})\ge1.7\times10^{-6}$--$7.5\times10^{-3}$
  \ \footnotesize(T109; $10^{2}$--$10^{6}$ below $\mu/2$ --- the cascade ends here)};
\foreach \a/\b in {rh/hand,hand/schur,schur/bare,bare/par,par/loew,loew/ratio,ratio/eps,eps/bnd,bnd/marg}
  \draw[ar] (\a) -- (\b);
\node[lbl] at (5.6,-0.78) {equivalence \emph{typing}};
\node[lbl] at (5.6,-2.33) {architecture of a};
\node[lbl] at (5.6,-2.63) {\emph{sufficient} route,};
\node[lbl] at (5.6,-2.93) {on finite windows};
\node[lbl] at (5.6,-8.53) {exact reduction};
\node[lbl] at (5.6,-10.08) {exact reduction};
\node[lbl] at (5.6,-11.63) {two certificates};
\node[lbl] at (5.6,-13.23) {one input left};
\node[draw=tfred!60!black,dashed,fill=tfred!4,rounded corners=2pt,inner sep=3.5pt,
      align=left,text width=46mm,font=\scriptsize,anchor=west] at (5.55,-5.4)
  {\textbf{The asymptotic wall.} Every step below the top line is established on
  \emph{16 resolved zones} at finite resolution. Sixteen zones are not all zones; no
  uniform-in-$k$ constant is certified anywhere in this cascade.};
\end{tikzpicture}
\caption{The reduction cascade, with its arrow types made explicit. The top line is an
\emph{equivalence typing} produced by an exact identity (Section~\ref{sec:i5}); the arrows below are the
architecture of a would-be sufficient route, verified on finitely many windows only. Since T109 the
cascade ends at one strict-positivity margin input; T110 shows that this input propagates itself along
the measured ladder (certified base case, graded Loewner minorant), leaving three gaps --- which T111
then drives deep and measures as three separate walls (the two subsections below). The dashed box is
the reason this figure is not a
proof sketch.}\label{fig:cascade}
\end{figure}

\subsection{The margin propagates: the circle closes on the measured zones (T110)}
T110 (\texttt{margin\_propagation\_probe.py}, contract \texttt{MARGIN.PROPAGATION}, $28/28$,
\textsc{margin-propagates} --- with a precision statement) asked the only question T109 left open:
does the strict scalar margin $m_k=\lambda_{\min}(Q|_{\mathrm{odd}})$ propagate itself along the zone
ladder, against the explicit demand $\mathrm{need109}(k)$ of the T109 chain? Three findings close the
circle on the measured zones; three gaps state precisely what a theorem would still need.

\emph{The margin map.} On a common grid ($16$ zones, $\delta_0=0.00284$) the margin falls
monotonically, $6.95\times10^{-2}$ ($n=2$) to $1.66\times10^{-5}$ ($n=29$), against
$\mathrm{need109}=3.09\times10^{-7}$--$3.87\times10^{-4}$: $m_k\ge\mathrm{need109}$ on $16/16$ zones,
with ratio $2.09$--$179.6$. Dilution during pure window growth is weak --- a factor $1.02$--$2.05$ per
atom gap, i.e.\ $0.993$--$0.997$ per cell, a $1-O(1/M)$ loss; the scaling fits are honestly poor
(the power law wins $7/16$). The structural finding of the part is the atom entry: it costs
\emph{nothing}. The naive expectation was a $-(\mu/2)$ pressure on the odd channel; measured, the
entry \emph{raises} $\lambda_{\min}$ on $15/15$ handovers (the first-order Kato term has helping
sign), and at $n=29$ the atom-free window is not even positive ($-9.2\times10^{-3}$): the atom makes
the window co-positive. The T106 killer argument --- transport destroys a demand that \emph{points}
somewhere --- has no analogue for the scalar floor.

\emph{The step law.} The growth step splits exactly (embedding error $\le2.6\times10^{-14}$), and
$\max|\mu N|=0.0$ on $15/15$: the new atom lies outside the old lag reach (gap $0.0606$ against
$\delta_0=0.00284$), so its restriction to the old window is the exact zero matrix --- everything is
decided by the bordering. Scalar routes fail decisively: bordered Weyl closes $0/15$ and the
Schur/Friedrichs scalar cap is vacuous $15/15$ (both recorded in the kill list). What passes certified
is the \textbf{graded Loewner minorant} --- its validity condition is exactly the induction
hypothesis, and the level is then found by bisection through a Cholesky/Sylvester factorisation: with
one soft direction ($n_{\mathrm{soft}}=1$) it holds $15/15$ with retention $1.0000$; the strictly
scalar variant ($n_{\mathrm{soft}}=\dim$) closes $0/15$.

\emph{Base case and circle test.} The base case $n=2$ is \emph{certified} by an explicit Cholesky of
$Q|_{\mathrm{odd}}-\lambda I$ at three grid resolutions ($\lambda_{\min}\ge6.93\times10^{-2}$, a
factor $88$--$180$ over need). The circle test then runs the whole chain from that certified base
value: the strictly scalar induction breaks at $k=3$ (its Schur factor compounds to
$2.7\times10^{-36}$), but the graded chain runs through completely --- $16/16$ zones above
$\mathrm{need109}$, each step one Cholesky, with the propagated floor as the only input. On the
measured zones, the induction circle closes end to end.

\begin{keybox}[The three gaps --- the asymptotic mountain, made precise]
\begin{enumerate}[leftmargin=1.6em]
\item \textbf{No reserve.} $f_{\mathrm{crit}}$ at the first handover is $1.00$; the chain lives on a
retention of $1-2\times10^{-7}$, and a factor-$2$ loss anywhere would break it within two steps. What
a theorem needs is a step law \emph{without} loss.
\item \textbf{No scalar step law.} Structurally excluded, not merely missed: the new cells attach at
the boundary layer where the pole source sits, so any one-number summary of the bordering is blind to
exactly the direction that decides it.
\item \textbf{No uniformity in $k$.} Every certificate above is one finite Cholesky; flat-in-$k$ is a
measurement, not a theorem. And the measured trend warns rather than reassures: $m_k\sim n^{-1.93}$
against $\mathrm{need109}\sim n^{-0.98}$, ratio $\sim n^{-0.96}$ (fit, rms $0.735$) --- extrapolated
crossing at $n\approx170$. Propagation gets \emph{harder} with $k$, not easier. (T111 subsequently
\emph{measured} that crossing: it sits at $n^*\approx462$, not $170$ --- the extrapolation was a fit
artefact of the $16$-zone window, but the wall itself is real. See the next subsection.)
\end{enumerate}
\end{keybox}

\subsection{The three walls (T111)}
T111 (\texttt{deep\_zone\_stress\_probe.py}, contract \texttt{DEEP.ZONE.STRESS}, $23/23$,
\textsc{crossing-confirmed}) stopped extrapolating and measured. On exactly the T110 cell grid ---
the first $16$ zones reproduce T110 bit-exactly --- the ladder was extended to \emph{every}
prime-power zone $n\le1000$ plus a thin tail to $3001$: $199$ zones in the margin map, with the
chain ladder running $n=2\ldots521$ through $117$ handovers.

\emph{The margin wall, measured.} The ratio $m_k/\mathrm{need109}$ falls $179.59\to0.0174$ across
the deep ladder, and the crossing is now a measurement, not a fit: the last zone with ratio $\ge1$
is $n=461$, the first below is $n=463$ (ratio $0.935$). The $n\approx170$ figure of T110 was a fit
artefact of the $16$-zone window --- the wall sits much further out, but it is real. The failure
follows the \emph{prime branch}: prime-power zones ($512$, $529$, $625$, $729$), whose atoms are
weaker, survive longer than the primes around them. The fit families were again compared honestly
(best by AIC: a broken-scale law with $n^*=501$; the jackknife band of the pure power law is
$[298,352]$), and a depth-preserving resolution study shows that refinement pushes the ratio
\emph{down}: $n^*\approx462$ is an \emph{upper} bound, not an estimate to be rounded up.

\emph{The ladder wall, arithmetic.} The uniformity object is no longer flat: $n^*_{\mathrm{soft}}$
runs $1\ldots8$ across the $117$ handovers. But retention stays $1.000000$ at every step (largest
single loss $5.96\times10^{-8}$), and the atom entry stays constructively free on $117/117$ --- the
restriction of the new atom to the old window is the exact zero matrix, deep into the ladder. The
new and harder finding is the \emph{ladder wall}: the nested step needs a lag gap exceeding twice
the window depth, and the twin-prime pair $521\to523$ (gap $0.003831<2D=0.005617$) kills the ladder
\emph{arithmetically} --- at frozen cell width, $n=521$ is the deepest reachable zone, and the bound
is set by twin primes, not by spectra. An arithmetic scan to $40\,000$ finds $4052/4284$
consecutive prime-power pairs endangered; letting $D$ shrink like $1/n$ instead forces the cell
count to grow like $n\log n/4$ (horizon $n\approx690$ at $h\le1500$).

\emph{The circle tears at the ratio, not at a step.} The no-reserve finding of T110 softens:
$f_{\mathrm{crit}}\sim n^{-0.39}$ (fit) --- the reserve \emph{opens} with depth, and the bottleneck
stays at the first step $2\to3$. The circle test runs from the certified base case (factor $179.6$)
and breaks first at $k=109$, $n=463$ --- exactly where the margin map measures the crossing. All
$117$ handovers certify at retention $1.000000$, including beyond the break: what fails is the
demand ratio, never the handover mechanism.

\begin{keybox}[The three walls --- and the reinterpretation]
The hardness balance names the driver: the growth of $1/\kappa$ (exponent $+0.86$) against the
margin (exponent $-0.68$). Three separate walls, in order of appearance:
\begin{enumerate}[leftmargin=1.6em]
\item \textbf{The margin wall}, $n^*\approx462$: measured, an upper bound (refinement pushes it
down); the T109-chain demand ratio crosses $1$ between $n=461$ and $n=463$.
\item \textbf{The ladder wall}, $521\to523$: purely arithmetic --- at frozen cell width the nested
handover dies at the first twin-prime gap below $2D$, independently of any spectral quantity.
\item \textbf{The requirement wall}, $n=727$: $\omega_{\mathrm{cert}}\ge1$ makes $\mathrm{need109}$
vacuous on $46$ deep zones --- beyond it the T109 chain no longer states a demand at all.
\end{enumerate}
The reinterpretation that comes with the measurement: the operating variable is the window
\emph{depth}, not the zone index $n$ --- doubling the cheeks moves the margin wall from $463$ to
$47$. A proof along this route must beat the exponent gap $-0.974=-0.681-0.293$ and couple the
depth $\delta_0$ to the local prime gap.
\end{keybox}

\noindent T112 subsequently ran the construction in exactly that coupled frame (next subsection):
walls (2) and (3) fall \emph{structurally}, wall (1) proves frame-invariant.

\subsection{The scaled frame: two walls fall, the margin wall is frame-invariant (T112)}
T112 (\texttt{adaptive\_scaling\_probe.py}, contract \texttt{ADAPTIVE.SCALING}, $20/20$,
\textsc{scaling-partial}) stopped freezing the cell width and coupled it to the local scale. Two
couplings were built: frame~A, gap-coupled, $D_k=g_k/(2\nu)$ with $g_k$ the local prime-power gap
and $\nu$ the resolution constant of the frame; and frame~B, mean-field, $D_k=\log n/(2\nu n)$ ---
the smooth PNT version of the same idea. The first finding is a structural asymmetry: \emph{frame~B
has no ladder} (only $38\%$ of consecutive pairs nest; the first failure is already at $n=3$) ---
the nested handover lives exclusively in the gap coupling, not in any smooth $1/n$ law. T105
admissibility, re-measured in the scaled frame, sets a resolution floor $\nu\ge4$ (at $\nu=1$,
$3/10$ test zones are inadmissible). The flatness verdict is honest and double-edged: the ratio is
\emph{not} flat, but the exponent is \emph{frame-invariant} --- the components shift massively
($m_k$: $n^{-1.93}\to n^{-0.95}$; $\mathrm{need109}$: $n^{-0.98}\to n^{+0.03}$, i.e.\ practically
$n$-independent in the scaled frame), yet the difference stays $-0.974$.

\emph{The three walls, scaled.} The \textbf{ladder wall is structurally gone}: with $D_k=g_k/(2\nu)$
the nested step has exactly $\nu$ cells per end \emph{by construction} --- twin pairs included ---
verified as an exact arithmetic lemma over all zones; the atom entry is the exact zero matrix by
construction rather than by arithmetic accident; and the T111 killer pair $461\to463$ now certifies
\emph{spectrally} ($h=1418$, retention $1.000000$), with the reserve \emph{opening}:
$f_{\mathrm{crit}}=8.1\times10^{-4}$--$4.2\times10^{-3}$ instead of $1.00$ --- roughly a factor
$1000$. (The pair $521\to523$ would need $h=1634$, above the dimension cap: arithmetically
co-certified, spectrally declared honestly as not run.) The \textbf{$\omega$ wall is no longer a
depth wall}: the depth coefficient is $n^{-0.079\pm0.121}$, compatible with zero ---
$\omega_{\mathrm{cert}}$ is driven by the resolution and the local gap, not by the depth of the
ladder; a residual requirement-vacuum tail of $5$ zones on the deepest ladder remains and is not
booked as ``gone''. The \textbf{margin wall stays}, practically unmoved: the first sub-$1$ zone
sits at $n=449$ against the frozen frame's $461$--$463$ (a factor $1.07$) --- it is not geometry
but the substance of the T109 requirement itself.

\emph{The limit object.} Overlaying the scaled windows, the archimedean part ($0.657$) and the pole
part ($0.558$) contract Cauchy-like with amplitude laws $D^{-0.26}$ and $D^{+0.37}$; the atom part
does \emph{not} contract ($0.935$) --- it is the prime-gap statistics itself, written in scaled
coordinates. The spectral shape does not drift: $\lambda_2/\lambda_1\sim n^{-0.029\pm0.052}$. The
resulting formulation is a split object: \emph{[a deterministic limit shape: arch $+$ rank-one
pole] $+$ [the atoms as a controlled, non-convergent perturbation]}. New and uncomfortable is the
\textbf{regrid object}: the zone frames are not nested, so Rayleigh--Ritz transfers nothing between
them; the regrid cost is measured at rate $(D'/D)^{+2.93}$ --- an object that simply does not exist
in the frozen frame, disclosed here because any limit argument has to pay it.

\begin{keybox}[The trade --- and the disclosed prime-gap dependence]
Three arithmetic walls are exchanged for two analytic demands:
\begin{enumerate}[leftmargin=1.6em]
\item[\textbf{[P1]}] \textbf{Positivity of the limit operator} --- \emph{one} operator instead of a
sequence of matrices; the deterministic part (arch $+$ rank-one pole) contracts, the atom part is
the perturbation to be controlled.
\item[\textbf{[P2]}] \textbf{A convergence rate below the step reserve} --- the Cauchy rates and the
regrid cost against the opening reserve $f_{\mathrm{crit}}$.
\end{enumerate}
The prime-gap inputs are disclosed rather than hidden. Upward, Bertrand--Chebyshev $1852$
($g_k\le\log2$; used and verified in-run) and the trivial even-gap bound suffice. The only remaining
arithmetic parameter is the normalised gap $\theta_k=g_k\,n/\log n=0.18$--$4.12$ --- and downward,
$\theta$ is \emph{provably} unboundedly small: by Zhang $2014$ / Maynard $2015$ (infinitely many
bounded gaps), $\theta\to0$ on subsequences, so any gap-coupled construction must be uniform in
$\theta\to0$, at the price of the $1/\theta$ cost explosion.
\end{keybox}

\noindent T113 subsequently answered the currency question (next subsection): the wall is
currency-invariant --- and what it measures is the discretisation, not the spectrum.

\subsection{The currency question: the wall is substance, and it measures the discretisation (T113)}
T113 (\texttt{limit\_operator\_probe.py}, contract \texttt{LIMIT.OPERATOR}, $27/27$,
\textsc{substance-confirmed}) answered the question T112 left open --- is the margin fall in the
scaled frame a normalisation artefact? --- and the answer reinterprets the program's hardest number.

\emph{The currency table.} The floor/requirement ratio carries the \emph{same} exponent,
$n^{-1.168\pm0.259}$, in all five currencies tested (raw, per $\lambda_{\max}$, per trace density,
per $D$, per $D^2$) --- and the invariance is structural, not accidental: the piecewise-constant
basis is $L^2$-orthonormal (Gram matrix exactly $I$, so there is no free $D$-power to spend), and
the zone matrices are restrictions of \emph{one} quadratic form to subspaces of the same $L^2$
space (cross-grid form identity at relative $1.2\times10^{-14}$). The floor's homogeneity degree in
the cell width is exactly $1.0$; the requirement's is $1.0$ under every legitimate rescaling. No
flat currency exists. The wall is currency-invariant: substance, not bookkeeping.

\emph{Two surprises about what that substance is.} First, the limit object is \emph{not} ``arch $+$
pole with the atoms as a perturbation'': the atom-free window form is deeply negative
($\lambda_{\min}=-3.08$ to $-19.36$, growing like $n^{+0.60}$), the atom block is equally large
(norm $3.21$--$19.44$), and the positive floor survives only as a \emph{cancellation} of relative
size $1.3\times10^{-7}$--$9.7\times10^{-5}$ --- on the floor vector, $-11.64$ against $+11.64$. The
Weyl bound is saturated: norm perturbation theory is five orders of magnitude too coarse for this
object. Second, there is \emph{no plateau under refinement}: at a fixed window, on nested grids
over $8$ levels (a factor $128$ in resolution), $\lambda_1\sim D^{1.83}$ and $\lambda_2\sim
D^{1.76}$ --- the \emph{same} power, with $\lambda_2/\lambda_1$ flat. \textbf{The continuum window
form has no gap; $m_k$ measures the gap of the discretisation.} This resolves the T112 tension
exactly: $\lambda_2/\lambda_1$ constant while the floor falls is what it looks like when both
eigenvalues carry the same $D$-power.

\emph{The atom part is almost invisible at the floor.} The entering (deepest) atom contributes
almost nothing to the coupling ($8.3\times10^{-11}$--$1.3\times10^{-6}$) --- which is why handover
steps certify at retention $1.000000$ while the margin tears. The floor is almost pure geometry:
$\log m=\mathrm{const}+1.87\log D-2.07\log\alpha$, with an arithmetic residual of a factor $1.13$
and no $\Lambda(n)$ dependence. The $\theta\to0$ stress test separates cost from substance: the
operator norm does \emph{not} diverge ($\theta^{-0.005}$); an extra atom at the window edge shifts
the floor by only $0.2$--$0.4\%$, while the same atom in the interior shifts it $10^7$ times more
--- the $1/\theta$ explosion of T112 is purely a \emph{cost} statement ($h\sim\nu n/\theta$ cells),
not an operator instability.

\emph{Regrid, re-read.} The raw regrid rate $(D'/D)^{3.85}$ is mostly the floor's own $D$-power ---
a normalisation jump, not an instability; the reserve certifies $f_{\mathrm{crit}}=9.5\times10^{-2}$
(better than the T112 estimate), and after subtracting the $D$-power the reserve buys $15$ steps
(against $0$ on the raw jump).

\begin{keybox}[The new hardness balance --- replacing {[P1]}/{[P2]}]
The T112 trade is restated by measurement; this balance supersedes the earlier two-demand form.
\begin{enumerate}[leftmargin=1.6em]
\item[\textbf{[P1]}] \textbf{Semidefiniteness of the balanced form without a gap} --- the continuum
window form is semidefinite at best; strict limit-floor formulations demand what refinement
demonstrably refuses.
\item[\textbf{[P2]}] \textbf{Certify the $D$-power itself} --- a Grenander--Szeg\H{o}-type
consistency statement for $\lambda_{\min}$ of the scaled form; the power, not an abstract rate, is
the object.
\item[\textbf{[P3]}] \textbf{New, and now leading: the requirement homogeneity} ---
$\mathrm{need109}$ carries $D^{0.63}$ against the floor's $D^{2.38}$: the T109 chain divides by an
\emph{artifact} margin. The repair the measurements point at is a \emph{margin-free step
certificate} --- a Schur/Loewner formulation certified from semidefiniteness alone, never dividing
by a margin.
\item[\textbf{[P4]}] \textbf{$\theta\to0$ is a cost-uniformity demand only} --- the frame needs
$h\sim\nu n/\theta$ cells, but the operator itself is $\theta$-stable.
\end{enumerate}
\emph{Disposition after T114 (next subsection):} [P3] is \textbf{discharged} --- the margin-free
step certificate exists and runs past the wall; [P4] stands as a pure cost statement; a new
\textbf{[P5]} appears (the (R) demand is grid-bound and is deleted, not transported); and
[P1]/[P2] are restated as the two-object core --- $\varepsilon$ relative to $\kappa$, and Schur
transport between grids.
\end{keybox}

\noindent T114 (\texttt{margin\_free\_step\_probe.py}, contract \texttt{MARGIN.FREE.STEP}) then
rebuilt the chain without the margin division --- and the wall dissolved.

\subsection{The wall dissolves: the margin-free step (T114)}
T114 (\texttt{margin\_free\_step\_probe.py}, contract \texttt{MARGIN.FREE.STEP}, $22/22$,
\textsc{wall-dissolves}) built the certificate that [P3] demanded, and the answer to the
preregistered question --- does the margin wall disappear as a requirement artefact? --- is yes.

\emph{The margin-free step.} The structural fact that makes it possible was already on the table
(T110/T111): the restriction of the grown window's form to the old window is bit-exactly
$Q_{\mathrm{old}}$ --- the entering atom's block is the exact zero matrix --- so margin-freeness
\emph{is} that property, once the step stops asking for more. Route (i) is Albert $1969$: the
generalised Schur criterion with the Moore--Penrose inverse,
\[
Q'=\begin{pmatrix}A & C\\ C^{\mathsf T} & X\end{pmatrix}\succeq0
\iff
X\succeq0,\quad \operatorname{ran}(C^{\mathsf T})\subseteq\operatorname{ran}(X),\quad
A-CX^{+}C^{\mathsf T}\succeq0
\]
(the range condition is Douglas $1966$) --- no margin appears anywhere. It certifies $27/27$
ladder steps, \emph{eleven of them beyond the old wall} ($467$, $479$, $512$, $529$, $547$,
$577$, $661$, $773$, $887$, $1129$, $1331$; deepest $h=1495$). The exact $4\times4$ Schur
complement is an $O(0.1)$ object with \emph{no} cancellation: $\lambda_{\min}(S)=0.068$--$0.154$,
i.e.\ $42$--$67\%$ of the block scale. Route (ii), the T110 graded minorant in its degenerate
limit $x_{\mathrm{in}}=0$, dies structurally on $27/27$ --- which retypes T110's margin: it was
not paying for positivity but for an \emph{abandoned direction}. Route (iii), unshifted Cholesky,
passes $27/27$ against the declared backward-error floor. Negative controls are sharp ($5/5$),
and the certificate is frame-invariant at $\nu=8$ ($3/3$). All seven zones at which the T109
chain tears --- including the wall zone $n=449$ itself, where $\mathrm{need109}/m=1.241$ --- are
certified by the margin-free step.

\emph{The wall in one line.} The same quantity, bounded through norms (Weyl:
$\lambda_{\min}(A)-\|C\|^2/\lambda_{\min}(X)$), is \emph{negative} by a factor of
$2.4\times10^{5}$--$9.6\times10^{7}$: an $O(1)$ numerator divided by the $10^{-6}$ artifact
floor. Every norm route had to fail --- and the exact Schur complement passes through none of it.

\emph{The favourable power direction.} On nested refinement chains at a fixed window,
$\varepsilon\sim D^{1.77\text{--}2.01}$ and $\kappa\sim D^{3.64\text{--}4.84}$ (fits) --- the
preregistered ``same power'' expectation is refuted \emph{in the favourable direction}: $\kappa$
falls faster, so the scale-free ratio $r=\kappa/\varepsilon\sim D^{1.6\text{--}3.1}$ shrinks
under refinement. The ratio closure $r=2\times10^{-4}$--$0.103$ holds on $27/27$ steps ($11/11$
beyond $463$). The circle [base semidefiniteness (a declared hypothesis input) $+$ margin-free
step $+$ ratio closure] closes $27/27$; ten multi-step chains run ($21$ certified steps, the
longest of length $4$); and \emph{every} chain ends at the cap $h\le1500$ --- a window
\emph{cost} --- never at a step. Regrid honesty: $0/464$ gap ratios are dyadic, so the
non-nested transport object of T112/T113 remains real.

\emph{[P5]: the (R) demand is grid-bound.} Holding the physical demand $(\mu/2)P_-$ fixed and
refining the grid, $\omega$ grows monotonically ($5/5$) and crosses $1$ ($4/5$): the (R)
formulation is \emph{false in the continuum} --- it is to be \emph{deleted}, not transported.
The exact Schur complement needs no demand surrogate in its place.

\emph{Cancellation robustness.} The $10^{-7}$-relative cancellation of T113 has $5.3$--$8.8$
decimal digits of headroom above the Cholesky backward-error floor (Wilkinson/Higham), and it
sits almost entirely in the rank-one pole --- whose positivity is exactly the T108 Szeg\H{o}
identity $\varepsilon$.

\begin{keybox}[The new core --- two objects]
The four-part balance collapses. [P3] is discharged (the margin-free step certificate exists and
runs past the wall), [P4] was only ever a cost statement, and [P5] deletes (R). What is left:
\begin{enumerate}[leftmargin=1.6em]
\item[\textbf{[P1]}] \textbf{$\varepsilon>0$ with controlled size relative to $\kappa$} --- the
exact Szeg\H{o}/Levinson identity of T108, now the single positivity object of the program.
\item[\textbf{[P2]}] \textbf{Transport of the exact Schur complement between non-nested grids}
--- an $O(0.1)$-sized object with no cancellation, sharpened from the earlier abstract
$D$-power demand.
\end{enumerate}
\emph{Disposition after T115 (next subsection):} [P2] \textbf{splits} --- the transport across a
real regrid certifies only mild refinement ($\rho\lesssim1.8$), for a principled reason (the
Schur floor itself falls under refinement), while a two-scale compression breaks the compute cap;
[P1] \textbf{sharpens} to one scalar inequality, the classical Szeg\H{o}--Levinson
prediction-error bound. The core is restated as the three-point list at the end of the next
subsection --- only one point of which is an inequality.
\end{keybox}

\noindent T115 (\texttt{schur\_transport\_probe.py}, contract \texttt{SCHUR.TRANSPORT}) then ran
the two last bricks --- and split them.

\subsection{Transport blocked, cap broken: the two-scale compression (T115)}
T115 (\texttt{schur\_transport\_probe.py}, contract \texttt{SCHUR.TRANSPORT}, $26/26$,
\textsc{transport-blocked} --- with a large compression gain) took the two-object core apart:
transport of the exact Schur complement across the ladder's non-nested grids, and a
multi-resolution compression of the already-certified interior.

\emph{The transport splits at $\rho^*\approx1.8$ --- for a principled reason.} Haynsworth's
partial minimisation $y^{\mathsf T}Sy=\min_z\,[y;z]^{\mathsf T}Q'[y;z]$ makes $\lambda_{\min}(S)$
available without any inverse and without a margin; the two-sided transport bracket built from it
is an \emph{identity} (bookkeeping), with the evidence sitting in the consistency defect $\eta$
and in the lower end. On the eleven real regrid pairs of the ladder ($\rho=1.023$--$3.425$) the
lower end is positive on $7/11$, with a clean split: every pair with $\rho\le2.061$ certifies and
none with $\rho\ge2.291$ does; a synthetic scan puts the threshold at $\rho^*=1.83$, against a
median real refinement ratio of $2.001$ --- $39\%$ of the ladder's $7686$ refinement pairs are
covered. The blocker is \emph{not} the transport machinery: the exact cell-overlap projection
beats three deliberately wrong transfer maps on $9/9$ (by factors up to $270$), and on
\emph{nested} ladders --- where the transport error is exactly zero --- $\lambda_{\min}(S)$
itself falls like $\rho^{-1.73}$ to $\rho^{-1.65}$. The drop is real, so no bound can undo it.
Two sharpenings came free: non-nestability is a matter of \emph{integrality}, not dyadicity
($0/14\,807$ gap ratios are integer, the closest miss $3.4\times10^{-5}$), and the measured
defect $\eta=2.7\times10^{-2}$--$8.1\times10^{-2}$ is $10^{3}$--$10^{6}$ times smaller than its
a-priori C\'ea/Strang norm surrogate --- which is exactly where the margin division would have
re-entered.

\emph{The compression breaks the cap.} The two-scale window (interior merged into blocks, the
boundary strip kept fine, the merge anchored at the centre) preserves
$X_{\mathrm{mixed}}=Q_{\mathrm{old,mixed}}$ \emph{exactly} (rel $0.0$) --- the compressed step is
still margin-free --- and Albert certifies $66/66$ (zone,\,$q$) combinations up to $q=64$, with
the compression error certified \emph{one-sided and second order} in the projection defect
(Rayleigh--Ritz plus stationarity of the fine minimiser; C\'ea/Strang; $66/66$) and savings down
to $m/h=0.043$. The deep run then breaks the T114 cap: the margin-free step certifies at
$n=155\,921$ ($117\times$ T114's $1331$), on a fine lattice of $h=93\,470$ cells ($62\times$ the
hard cap $h\le1500$) compressed to $m=1490$; re-coarsening inside a chain is free (Rayleigh--Ritz,
rel $\le2.1\times10^{-12}$); the longest chain is now \emph{ten} steps (T114: four), with $33$
certified steps over four chains; the stopper is the cost cap on $3/4$ runs and a failing step on
$0$; and every certified quantity sits $10^{5}$--$10^{11}$ times above the declared Cholesky
backward-error floor. Stated honestly: $\lambda_{\min}(S_{\mathrm{mixed}})\approx5.1$ is $52\times$
the uniform value because the coarser Galerkin space sees fewer soft directions --- a weaker
statement about the continuum, declared as such, and flat across the whole deep band. The
contiguous reach an induction needs moves from $n\le125$ to $n\le5437$ (factor $43.5$): the
$n\log n$ cost barrier is \emph{divided}, not broken.

\begin{keybox}[The remaining list --- three points, one inequality]
The shortest list the program has produced. From T114's balance, the margin (dissolved), the (R)
demand (grid-bound, deleted), the floor/$\varepsilon$ transport (the object to transport is
$\lambda_{\min}(S)$, the slowest-falling of the three) and the compression risk (it does not
break the step) are all \emph{struck}. What is left:
\begin{enumerate}[leftmargin=1.6em]
\item[\textbf{(1)}] \textbf{[P1] as one scalar inequality} --- given
$T=Q+\tilde t\tilde t^{\mathsf T}\succ0$, positivity of $Q|_{\mathrm{odd}}$ is \emph{equivalent}
to $\varepsilon=1-\tilde t^{\mathsf T}T^{-1}\tilde t\ge c(D)>0$: the classical
Szeg\H{o}--Levinson prediction-error bound for one explicit symbol (the Haynsworth double-Albert
device: $T\succ0\wedge\varepsilon\ge0\iff Q\succeq0$). This is the target-sentence candidate.
Measured: $\varepsilon\sim\rho^{-1.71}$, falling slightly faster than
$\lambda_{\min}(S)\sim\rho^{-1.69}$ (fits).
\item[\textbf{(2)}] \textbf{An a-priori bound on the transport defect $\eta$} --- a regularity
statement about the minimiser $z^*$; it would turn the measured transport into a certificate for
$\rho\le1.8$.
\item[\textbf{(3)}] \textbf{A boundary formulation} (suggested by the exact-zero lemma) ---
never represent the old interior at all; it would remove the compute cap \emph{and} the regrid
at once.
\end{enumerate}
Only (1) is an inequality; (2) is a regularity statement, (3) is a reformulation.
\end{keybox}

\noindent T116 (\texttt{boundary\_formulation\_probe.py}, contract \texttt{BOUNDARY.FORMULATION}
--- the induction as a pure boundary process: a Riccati-type Schur recursion, the pole via
Woodbury, kernel-tail truncation; a deep boundary run; the $\varepsilon(D)$ law against the
classical benchmarks) was running at the time of writing and is \emph{not} reported here.

% =====================================================================
\section{The kill list}\label{sec:kills}
% =====================================================================
\noindent A program that only reports what survived is not reporting. These are the readings and routes
that were refuted \emph{by the program's own probes}, each with the measurement that killed it. They are
listed as a section rather than as footnotes because they are the more transferable output: each one
tells a later attempt where not to go.

\begin{enumerate}[leftmargin=1.6em,label=\textbf{K\arabic*.}]
\item \textbf{The essential-singularity reading of the onset} (T96, refuted T102). T96 read the
counterfactual onset as super-polynomial, compatible with $\exp(-c/\delta')$, i.e.\ an essential
singularity. T102 showed the onset is \emph{anchored} at the crossing $\delta_c=2w_k$ of two finite
quantities ($R^2=0.968$), that exponential versus power behaviour is not separable at this resolution,
and that the scatter of $c$ ($\mathrm{CV}$ $103\%$) collapses to $c'=1.52$ in the scale-free variable.
The essential-singularity reading is \emph{compatible but not distinguished} --- and it is no longer
used.
\item \textbf{The $C/g$ law as a causal law} (T101, refuted T102). Triple negative: $g_k$ is causally
impossible ($Q_{k-1}$ is blind to atom $k+1$), statistically dispensable (causal law $C\cdot u$:
$\mathrm{CV}$ $23.3\%$ versus $27.6\%$), and arithmetically only a ceiling ($C_k\le g_k\mu_k$ exactly;
the $16$-zone extrapolation violates its own ceiling from $k=69$; checked on $18\,120$ prime-power
atoms to $n=200\,000$). \emph{It was a proxy.}
\item \textbf{The naive margin route} (T104). $M\ge m\cdot\mathrm{Id}$ holds in $0/16$ zones; the
coupling-mass to margin ratio is $10^3$--$10^6$; $m\sim M^{-1.7}$ vanishes in the continuum (fit). The
genuine form is nearly singular on the induction window and no margin rescues it. (The prolate wing
basis recommended by T96/T103 was tested in both arms and also rejected.)
\item \textbf{The density chain} (T106). It would need a hard gap $w_0\ge0.55$--$2.44$; measured
$5.3\times10^{-6}$--$1.1\times10^{-3}$ --- short by a factor $5\times10^2$ to $3.6\times10^5$. Dead for
good.
\item \textbf{The Landau--Widom anchor} (T106). The wrong classical anchor for this problem
(rms $0.50$--$1.55$). The right one is Grenander--Szeg\H{o} on the Planck-coarsened symbol --- averaging
over the mode-spacing cell $\pi/M$ --- which hits $0.92$--$2.43$ and is better on $67/72$ pairs. This
also explains why a negative symbol (dip $-21.3\ldots-3.4$) still gives $Q\succeq0$: the atoms oscillate
exactly at the window's Planck scale.
\item \textbf{The accumulated invariant with self-propagation} (T106). A no-go \emph{with a measured
mechanism}. The wing demand gives $\beta_{\mathrm{wing}}=2.21$--$14.43>1$ on $16/16$, but inherited
demands give $\beta_{\mathrm{int}}\sim10^{-2}$ on $15/15$: an accumulated invariant with absolute
anchoring is false. The reason is transport --- in a larger window an old wing pair becomes an interior
pair and draws $99.7$--$99.95\%$ of its demand from the $64$ softest modes (coupling $751$--$2653\times$
stronger). No self-propagation either ($\beta_{\mathrm{on}}<1$ on $15/15$, loss $13$--$320\times$,
$\theta_k=0.9866$--$0.9998$). \emph{The wing margin is an edge property of the window, not a property of
the pair.} What survives is the window-wise family $\sigma_k(\delta_0)\ge1.31(\mu_k/2)$, uniform in $M$
with drift $\le2.01\times$.
\item \textbf{The symbol route for the odd channel} (T107). Structurally dead, not merely weak: the
symbol lower bound for $T_{\mathrm{odd}}$ is certified but \emph{negative}
($\lambda_{\mathrm{cert}}=-2.54\ldots-0.81$) against a measured
$\lambda_{\min}(T_{\mathrm{odd}})=4.7\times10^{-5}\ldots6.4\times10^{-2}$ and a symbol infimum of
$-22\ldots-6$: five orders of magnitude and a sign change. The soft edge of $T_{\mathrm{odd}}$ is a pure
finite-section effect that no symbol can see. Coarse certified chains (Cauchy--Schwarz against
$\lambda_{\min}$) close $0/16$ --- and fail even with the \emph{measured} $\lambda_{\min}$, so it is a
category error, not a precision problem.
\item \textbf{Three certificate candidates for the avoidance norm} (T108). Killed \emph{quantitatively},
which is the useful form: the cancellation route ($\sum|k|=401$ against $\sum k=1$), the Bessel route
(off by $10^7$) and the pole Cauchy--Schwarz route (off by $10^6$). $\omega_{\mathrm{cert}}$ is blocked
by the compression Schur distance alone ($0.18$--$0.52$). What replaced them is not a fourth bound of
the same kind but a different object class: a boundary-decay statement about one explicit vector.
T109 then killed the natural fourth candidate as well: \textbf{the Combes--Thomas resolvent bound at
the boundary}, evaluated with the \emph{exact} Green row so that no constant can rescue it, closes
$1/16$ --- the mechanism at the edge is cancellation ($44$--$4.6\times10^{4}$ on $15/16$ zones), which
no absolute-value bound survives. What works instead is the residual certificate that carries the
cancellation (Section~\ref{sec:certification}).
\item \textbf{The scalar step law for the margin} (T110). Killed structurally, in three pieces:
bordered Weyl closes $0/15$ handovers; the Schur/Friedrichs scalar cap is vacuous on $15/15$ (its
scalar-only bound degrades back to Weyl); and the strictly scalar graded minorant
($n_{\mathrm{soft}}=\dim$) closes $0/15$, with the strictly scalar \emph{induction} compounding to a
Schur factor of $2.7\times10^{-36}$ by $k=3$. The reason is a boundary layer: the new cells of a
grown window attach at the edge where the pole source sits, so a one-number summary of the bordering
is blind to the deciding direction. What survives is the graded ($n_{\mathrm{soft}}=1$) Loewner
minorant --- one soft direction is exactly the amount of geometry the step needs.
\item \textbf{Two earlier structural kills}, recorded here because they still bound the design space:
the cross-block $\sqrt{\lambda_0}$ law is Douglas range inclusion and therefore circular (T98), and the
cone-library route saturates at $5/24$ coverage and is dead as a covering strategy (T72, promoted into
\vref{v540} as a named limit). Likewise the seam route to the archimedean place (T20--T25) and the
Kohnen-plus route to the central values (T44, promoted into \vref{v537} as a scope fence).
\end{enumerate}

% =====================================================================
\section{Status and outlook}\label{sec:status}
% =====================================================================
\subsection{What is certified, what is measured, what is hypothesis input}
\begin{keybox}[The three-way split, stated once]
\textbf{Certified} (exact identities or closed-form bounds, machine-checked): the seven ledger modules
of Section~\ref{sec:verified}; and, in the sandbox arc, the bare bound $b_0$ in closed form, the parity
superselection $JQJ=Q$, $JB=-BR$, the pole splitting
$ab^{\mathsf T}+ba^{\mathsf T}=ss^{\mathsf T}-tt^{\mathsf T}$, the channel inertia, the avoidance law
$\Qfull\succeq0\Rightarrow\|B^{\mathsf T}v_i\|^2\le w_i\Lambda_-$, the growth inequality, the
accumulation bounds, the Sherman--Morrison/Woodbury reduction of (R) to $r\le1$, the three T108
identities ($\varepsilon$ as a $Q|_{\mathrm{odd}}$ energy, the rank-one overlap $hh^{\mathsf T}/
\varepsilon$, and $\varepsilon$ as the square of the last Cholesky pivot), the two T109
certificates (the graded matrix cap $\omega_{\mathrm{cert}}=0.2561$--$0.7569$, no hypothesis input,
and the residual certificate for the boundary value, sharpness $1.0000$--$2.4937$, $16/16$ in
budget), the three T110 pieces of the propagation circle (the exact splitting of the growth step,
embedding error $\le2.6\times10^{-14}$, with the new atom's restriction to the old window the exact
zero matrix; the graded Loewner minorant step, valid exactly under the induction hypothesis, holding
$15/15$ with retention $1.0000$; and the base-case Cholesky certificate at $n=2$,
$\lambda_{\min}\ge6.93\times10^{-2}$ at three resolutions), and their T111 deep-ladder extension
($117$ handover certificates at retention $1.000000$, largest single loss $5.96\times10^{-8}$; atom
entry the exact zero matrix on $117/117$; and the arithmetic ladder-wall statement at frozen cell
width, gap $>2D$ needed against the twin-prime gap $0.003831$ at $521\to523$), the T114
margin-free step certificates (Albert/pseudo-inverse Schur, $27/27$ ladder steps to $n=1331$
including all seven torn zones, with the exact $4\times4$ Schur complement at $42$--$67\%$ of the
block scale --- no margin division anywhere), and the T115 compression certificates (the exact
identity $X_{\mathrm{mixed}}=Q_{\mathrm{old,mixed}}$; Albert on $66/66$ (zone,\,$q$) combinations;
the one-sided second-order compression error; the deep margin-free step at $n=155\,921$ and the
ten-step chains --- certificates for a coarser Galerkin space, declared as such).\\[3pt]
\textbf{Measured} (holds on the resolved windows; no uniform constant): $\rho_-$; the ratio
$r=0.0053$--$0.1814$ and $\hat r<1$ at all $92$ resolution points; $\omega=0.2366$--$0.7531$; the
strict odd-channel margin $\lambda_{\min}(Q|_{\mathrm{odd}})=1.4\times10^{-5}$--$5.1\times10^{-2}$; the
margin map $m_k\ge\mathrm{need109}$ on $16/16$ zones (ratio $2.09$--$179.6$), extended by T111 to
$199$ zones with the crossing \emph{measured} between $n=461$ and $n=463$ ($n^*\approx462$, an upper
bound; the T110 extrapolation $n\approx170$ was a fit artefact); the reserve trend
$f_{\mathrm{crit}}\sim n^{-0.39}$ (fit); the requirement wall $n=727$; the closure map $44/64$; the
handover margins; the two race slopes; the T114 ratio closure $r=2\times10^{-4}$--$0.103$ on
$27/27$ steps and the $D$-powers of $\varepsilon$ and $\kappa$ (fits); the T115 transport split
(the bracket's lower end positive on $7/11$ regrid pairs, threshold $\rho^*=1.83$ against a
median ratio $2.001$) and the fall laws $\lambda_{\min}(S)\sim\rho^{-1.69}$,
$\varepsilon\sim\rho^{-1.71}$ (fits); every scaling
exponent quoted in this document.\\[3pt]
\textbf{Hypothesis input / classical} (used, named, not proved here): Kneser, Witt/Eichler, Shimura,
Waldspurger/Baruch--Mao, Cohen, Weil--Guinand, Bombieri, Yoshida, Connes--Consani, Douglas,
Slepian--Landau--Pollak, Gronwall/Robin (unconditional form only), Grenander--Szeg\H{o},
Sherman--Morrison/Woodbury, Christoffel--Darboux, Szeg\H{o}--Levinson, the Weyl bound,
Combes--Thomas/Jaffard, Prager--Synge, first-order Kato perturbation, Albert (the pseudo-inverse
Schur criterion), Haynsworth (the partial-minimisation form of the Schur complement),
C\'ea/Strang's first lemma, Bramble--Hilbert, Brandt/Yserentant two-scale spaces,
Wilkinson/Higham backward error.
\end{keybox}

\subsection{The asymptotic mountain}
Three separate reasons why the cascade of Figure~\ref{fig:cascade} is not a route to a theorem as it
stands, and why the document does not present it as one. Since T110 the mountain has a
\emph{precise} formulation --- the three gaps of Section~\ref{sec:certification}: no reserve, no
scalar step law (structurally excluded by the boundary layer), and no uniformity in $k$. Since T111
it also has a \emph{measured} one --- the three walls: the margin wall at $n^*\approx462$ (measured,
an upper bound; the earlier $n\approx170$ was a fit artefact), the arithmetic ladder wall at the
twin-prime pair $521\to523$, and the requirement wall at $n=727$. The mechanism itself never failed
on the deep ladder --- all $117$ handovers certify at retention $1.000000$, and the reserve even
opens with depth ($f_{\mathrm{crit}}\sim n^{-0.39}$, fit) --- so the actionable content of the
measurement is the reinterpretation: the operating variable is the window depth, not the zone index,
and a proof along this route must beat the exponent gap $-0.974=-0.681-0.293$ and couple the depth
$\delta_0$ to the local prime gap. T112 (\texttt{adaptive\_scaling\_probe.py}) then performed that
coupling: the ladder wall and the depth dependence of the requirement wall fall \emph{structurally}
in the gap-coupled frame, but the margin wall is frame-invariant at exponent $-0.974$, with the
prime-gap dependence disclosed (Bertrand--Chebyshev suffices upward; Zhang/Maynard force
$\theta$-uniformity downward). T113 (\texttt{limit\_operator\_probe.py}) then settled what that
invariant number \emph{is}: the wall carries the same exponent in every admissible currency --- it
is substance --- but it measures the \emph{discretisation}, not the spectrum: at a fixed window
both leading eigenvalues carry the same cell-width power under refinement (the continuum window
form has no gap), the positive floor is a cancellation of relative size $10^{-7}$--$10^{-4}$, and
the T109 requirement chain divides by an artifact margin. T114
(\texttt{margin\_free\_step\_probe.py}) then discharged the leading part of that balance: the
margin-free step certificate exists (Albert $1969$, the exact Schur complement), runs eleven
steps past the old wall to $n=1331$, certifies all seven torn zones, and shows the wall was an
$O(1)$ numerator divided by an artifact floor --- chains now stop only at the $h\le1500$ window
cost, never at a step, and the (R) demand itself is grid-bound and deleted ([P5]). T115
(\texttt{schur\_transport\_probe.py}) then split the two-object core: the transport of the exact
Schur complement across a real regrid certifies only mild refinement ($\rho\lesssim1.8$) --- and
for a principled reason, since on nested ladders, where the transport error is exactly zero,
$\lambda_{\min}(S)$ itself falls like $\rho^{-1.7}$ --- while the two-scale compression breaks
the compute cap: a certified margin-free step at $n=155\,921$ ($117\times$ deeper), ten-step
chains, the stopper always cost and never a step, with the contiguous reach moving from
$n\le125$ to $n\le5437$ (the $n\log n$ barrier divided, not broken). Since T115 the mountain is
the three-point list of Section~\ref{sec:certification} --- an a-priori bound on the transport
defect, a boundary formulation, and exactly \emph{one} inequality: the classical
Szeg\H{o}--Levinson prediction-error bound $\varepsilon\ge c(D)>0$ for one explicit symbol. T116
(\texttt{boundary\_formulation\_probe.py}) was running at the time of writing.

\begin{enumerate}[leftmargin=1.6em]
\item \textbf{Sixteen zones are sixteen zones.} Everything below the top line of the cascade is
established on the resolved zones $n=2\ldots29$ (T111's deep ladder extends the measurements to
$n=521$, T114's margin-free ladder to $n=1331$, T115's compressed run to $n=155\,921$ --- finitely many zones either way). Not one constant in the chain is certified uniformly in
$k$. The T101 requirement list (A)--(D) is unchanged: the hardness would sit in (A), a sharp,
non-perturbative arithmetic lower bound at a Weil-positivity edge.
\item \textbf{Finite resolution is a wall, and a measured one.} Numerics as a witness is exhausted
beyond $\alpha\approx0.55$ ($O(M^{-1.8})$, T96); the finite-block route was judged the structurally
wrong instrument, not an under-resourced one (T92); $\lambda_{\min}$ collapses geometrically by a factor
$11.4$ per mode, so double precision ends at $N\approx10$ while $N=24$ would need $10^{-27}$ entries.
\item \textbf{The induction base is not free either.} The recursion terminates arithmetically
($240/240$ in $\le4$ steps to the classical zone, T99; $960/960$ with longest chain $13$ steps, T101),
but the zone peaks need percent-level bounds, not $0.1\%$ (T99: extrapolated peak margins
$+2.7\times10^{-2}$, $+1.4\times10^{-2}$, $+1.7\times10^{-2}$, i.e.\ $2.2$--$5.5\%$ of $\mu_k/2$, with
zone $5$ saturating exactly).
\end{enumerate}

\subsection{Extraction candidates: the community-checkable path}
The useful question at this point is not ``how close is the program to RH'' --- it is not close, and no
such claim is made --- but ``which pieces of it are unconditional theorems that someone else can
check''. Two candidates are named in the diary and are repeated here as the outlook:

\begin{itemize}[leftmargin=1.5em]
\item \textbf{Zone extension beyond $\log2$.} Unconditional Weil positivity is classically known exactly
up to support width $\log2$ (Yoshida 1992; Bombieri 2000). T91's target inequality --- for
$a\in(\log2,\,a^*]$ and $\|f\|=1$, \ $(P_{\mathrm{pole}}+A_{\mathrm{arch}})(f)\ge\sqrt2\log2\cdot h_f(\log2)$
--- is typed as \emph{provable-shaped} and would, if proved, be a standalone classical result: a zone
extension past Bombieri's boundary. Its typing is explicit and unflattering: RH $\Rightarrow$ (T), but
(T) $\not\Rightarrow$ RH. T92 showed the finite-block route to it is the wrong instrument; that
narrows, rather than closes, the target.
\item \textbf{The matching lemma.} Already proved exact-integer on $[4,10^6]$ with an explicit margin
inside \vref{v541}, and reduced by T77 to a restricted $\sigma_{-1}$ divisor bound in the Gronwall
regime --- with the note that Robin's \emph{unconditional} 1983 bound suffices at the measured
headroom, so the RH-equivalent sharpening is not needed. Closing the tail (the correlated-cancellation
lemma on credit-rich non-coherent atoms) would be a self-contained classical statement.
\end{itemize}

\begin{gapboxnb}[Final fence]
The prime line of TFPT has produced seven machine-checked modules about lattice-native Hecke structure,
a finite relative-trace identity, and an exactly closing transport ledger. It has \emph{not} produced
evidence for the Riemann Hypothesis, and the one place where RH appears --- I5 --- appears as an
equivalence typing that names the residual difficulty rather than reducing it. The corresponding open
gates in the ledger (\texttt{ZETA.CENSUS.TO.GL1}, \texttt{ZETA.HP.CARRIER}) are untouched by everything
in Sections~\ref{sec:i5}--\ref{sec:certification}, and no marker moves anywhere on account of this
document.
\end{gapboxnb}

% =====================================================================
\appendix
\section{Probe index, diary parts T11--T115}\label{app:index}
% =====================================================================
\noindent Every row is one exploratory probe in \texttt{experiments/tfpt-discovery/}; the file column
omits the common suffix \texttt{\_probe.py}. ``Checks'' is the probe's own assertion count, all green.
Verdict strings are the diary's preregistered verdict enums. Rows marked ``$\to$ \vref{vN}'' were later
consolidated into the corresponding load-bearing module; the promotion, not the probe, is the
load-bearing artefact. The index covers the $105$ parts T11--T115 that the current diary file records
($106$ runs --- T104 was a double run with two independent arms --- and $2684$ green checks); the ten
earlier probes T1--T10 predate that file and are not indexed here. Series totals as recorded by the
diary at T115: \textbf{$115$ probes, $2876/2876$ sandbox checks}, plus \vref{v535}--\vref{v541}
promoted. T116
(\texttt{boundary\_formulation\_probe.py}, contract \texttt{BOUNDARY.FORMULATION}) was still running
when this document was written and is deliberately absent.

{\footnotesize
\setlength{\tabcolsep}{3.2pt}
\begin{xltabular}{\linewidth}{@{}p{0.95cm} >{\raggedright\arraybackslash}p{3.1cm} Y >{\raggedright\scriptsize\arraybackslash}p{3.35cm} p{1.15cm}@{}}
\toprule
\textbf{Part} & \textbf{Probe} & \textbf{Contract / finding} & \textbf{Verdict} & \textbf{Checks}\\
\midrule
\endfirsthead
\toprule
\textbf{Part} & \textbf{Probe} & \textbf{Contract / finding} & \textbf{Verdict} & \textbf{Checks}\\
\midrule
\endhead
\bottomrule
\endfoot
T11 & \texttt{\scriptsize e8\_\allowbreak glue\_\allowbreak chi4\_\allowbreak signed\_\allowbreak theta} & the $\chi_4$-signed glue theta welds parts 6 and 7 & NO PROMOTION & 19/19\\
T12 & \texttt{\scriptsize e8\_\allowbreak glue\_\allowbreak lseries\_\allowbreak selfdual} & L-series of the signed glue count; Ap\'ery/$\zeta(3)$ bridge; Poisson self-dual width $2\pi$ & NO PROMOTION & 16/16\\
T13 & \texttt{\scriptsize e8\_\allowbreak glue\_\allowbreak character\_\allowbreak atlas} & glue-character atlas of $E_8$ & NO PROMOTION & 11/11\\
T14 & \texttt{\scriptsize seam\_\allowbreak selfdual\_\allowbreak width} & SEAM.SELFDUAL.WIDTH; unimodularity gate & DEFLATED & 13/13\\
T15 & \texttt{\scriptsize zeta\_\allowbreak arith\_\allowbreak completion} & ZETA.ARITH.COMPLETION; shell Hamiltonian & PARTIAL & 12/12\\
T16 & \texttt{\scriptsize zeta\_\allowbreak local\_\allowbreak euler} & ZETA.LOCAL.EULER; one weight-4 Hecke polynomial & PARTIAL & 16/16\\
T17 & \texttt{\scriptsize zeta\_\allowbreak funceq\_\allowbreak duality} & ZETA.FUNCEQ.DUALITY; Poisson involution & PARTIAL & 9/9\\
T18 & \texttt{\scriptsize zeta\_\allowbreak weil\_\allowbreak recovery\_\allowbreak trace} & ZETA.TRACE.COMPLETENESS; Guinand--Weil balance & BENCHMARK-STANDING & 17/17\\
T19 & \texttt{\scriptsize prime\_\allowbreak thinning\_\allowbreak seam\_\allowbreak truncation} & PRIMON.SEAM.TRUNCATION; thinning vs.\ the finite tower & MIXED & 19/19\\
T20 & \texttt{\scriptsize seam\_\allowbreak archimedean\_\allowbreak kernel} & archimedean kernel from the seam density of states & KILLED-AS-NAIVE & 25/25\\
T21 & \texttt{\scriptsize prime\_\allowbreak prediction\_\allowbreak mechanics} & four prediction channels quantified & MAPPED & 19/19\\
T22 & \texttt{\scriptsize seam\_\allowbreak halfcut\_\allowbreak dilation\_\allowbreak arch} & non-DOS successor of the T20 kill & PARTIAL (K1,\,K2 pass; K3 fails) & 25/25\\
T23 & \texttt{\scriptsize seam\_\allowbreak boost\_\allowbreak arch\_\allowbreak dictionary} & arch dictionary via the seam boost & DEAD & 23/23\\
T24 & \texttt{\scriptsize seam\_\allowbreak truncated\_\allowbreak trace\_\allowbreak rvm} & truncated trace / Riemann--von Mangoldt route & DEAD & 25/25\\
T25 & \texttt{\scriptsize seam\_\allowbreak tate\_\allowbreak phase} & scattering phase as the last observable & DEAD (arch route closed) & 29/29\\
T26 & \texttt{\scriptsize shell\_\allowbreak multiplicity\_\allowbreak one\_\allowbreak quotient} & shell quotient towards multiplicity one & PARTIAL & 18/18\\
T27 & \texttt{\scriptsize e8\_\allowbreak pneighbor\_\allowbreak hecke} & Kneser $p$-neighbours carry the Eisenstein half & PARTIAL $\to$ \vref{v535} & 27/27\\
T28 & \texttt{\scriptsize e8\_\allowbreak kneser\_\allowbreak transport\_\allowbreak holonomy} & cusp coefficient from mod-3 lattice geometry & TRANSPORT-VISIBLE & 22/22\\
T29 & \texttt{\scriptsize primon\_\allowbreak transport\_\allowbreak category} & primon monoid relations & RELATIONS-VERIFIED & 16/16\\
T30 & \texttt{\scriptsize e8\_\allowbreak transport\_\allowbreak cusp\_\allowbreak uniformity} & uniformity of the T28 reading at $p=5,7$ & BROKEN (reading); transport holds & 32/32\\
T31 & \texttt{\scriptsize e8\_\allowbreak neighbor\_\allowbreak operator\_\allowbreak decomposition} & cusp extraction is a rule $\nu_p=a\,\mathrm{Id}+b\,T_p$ & REPAIRED-BY-DICTIONARY $\to$ \vref{v535} & 49/49\\
T32 & \texttt{\scriptsize census\_\allowbreak newform\_\allowbreak recovery} & census redundancy is 2-adic oldform structure & REDUNDANCY-IS-OLDFORM $\to$ \vref{v535} & 24/24\\
T33 & \texttt{\scriptsize e8\_\allowbreak nu\_\allowbreak rule\_\allowbreak many\_\allowbreak primes} & hardening of the $\nu$-rule to $p\le100$; Eichler finding & HARDENED $\to$ \vref{v536} & 26/26\\
T34 & \texttt{\scriptsize compiler\_\allowbreak functor\_\allowbreak theta} & project ZETA.COMPILER.FUNCTOR founded & FOUNDED-WITH-ANCHOR & 13/13\\
T35 & \texttt{\scriptsize rankin\_\allowbreak selberg\_\allowbreak functor} & Rankin--Selberg translates the abelian shadow & RS-TRANSLATES-EISENSTEIN & 13/13\\
T36 & \texttt{\scriptsize e8\_\allowbreak lambda\_\allowbreak charsum\_\allowbreak evaluator} & cusp remainder isolated geometrically & TWO-SIDED $\to$ \vref{v536} & 29/29\\
T37 & \texttt{\scriptsize e8\_\allowbreak salie\_\allowbreak signed\_\allowbreak cusp} & signed cusp channel, scalable & SCALABLE-AND-SIGNED & 37/37\\
T38 & \texttt{\scriptsize cuspidal\_\allowbreak bridge} & Shimura preimage of $f_8$ inside the compiler monoid & BRIDGE-FOUND $\to$ \vref{v537} & 13/13\\
T39 & \texttt{\scriptsize zeta\_\allowbreak center\_\allowbreak atlas} & the abelian drop & ABELIAN-DROP-CLOSED & 13/13\\
T40 & \texttt{\scriptsize weil\_\allowbreak recovery\_\allowbreak compiler\_\allowbreak sector} & Weil-recovery sector terrain & TERRAIN-MAPPED & 24/24\\
T41 & \texttt{\scriptsize functor\_\allowbreak wiring\_\allowbreak halfintegral} & the half-integral bridge is Hecke-equivariant & FUNCTOR-WIRED $\to$ \vref{v537} & 13/13\\
T42 & \texttt{\scriptsize e8\_\allowbreak local\_\allowbreak densities} & local densities closed & CLOSED $\to$ \vref{v536} & 25/25\\
T43 & \texttt{\scriptsize kohnen\_\allowbreak plus\_\allowbreak waldspurger} & Kohnen plus space & PARTIAL & 12/12\\
T44 & \texttt{\scriptsize kohnen\_\allowbreak operator\_\allowbreak projection} & central-value route via Kohnen plus & OUTSIDE-KOHNEN-SCOPE $\to$ \vref{v537} & 17/17\\
T45 & \texttt{\scriptsize baruch\_\allowbreak mao\_\allowbreak metaplectic} & central-value family as $b(d)^2$; $R$ constant & CENTRAL-VALUES-WIRED $\to$ \vref{v537} & 13/13\\
T46 & \texttt{\scriptsize stage4\_\allowbreak prolate\_\allowbreak prototype} & stage-4 prolate prototype & IMPL-OPEN & 16/16\\
T47 & \texttt{\scriptsize stage4\_\allowbreak prolate\_\allowbreak uv\_\allowbreak dirac} & pointwise UV/Dirac prototype & UV-NO-SIGNAL & 17/17\\
T48 & \texttt{\scriptsize lambda\_\allowbreak scale\_\allowbreak uniqueness} & does the compiler fix the $\Lambda$ scale? & SCALE-TORSOR (claim killed) & 17/17\\
T49 & \texttt{\scriptsize rankin\_\allowbreak positivity\_\allowbreak miniature} & the Deligne/Rankin mechanism runs in-suite & MINIATURE-RUNS & 28/28\\
T50 & \texttt{\scriptsize half\_\allowbreak integral\_\allowbreak selfconvolution} & the central-value family aggregates exactly & FAMILY-AGGREGATED & 17/17\\
T51 & \texttt{\scriptsize waldspurger\_\allowbreak family\_\allowbreak kernel} & first non-collapsing infinite carrier candidate & INFINITE-CARRIER-CANDIDATE & 22/22\\
T52 & \texttt{\scriptsize hecke\_\allowbreak orbit\_\allowbreak transfer} & path space of Hecke words; fibre determinants & TRANSFER-REALIZED & 16/16\\
T53 & \texttt{\scriptsize relative\_\allowbreak trace\_\allowbreak compiler} & \vref{v535}/\vref{v536}/\vref{v537} are three projections of one RTF identity & ONE-FORMULA $\to$ \vref{v538} & 16/16\\
T54 & \texttt{\scriptsize waldspurger\_\allowbreak canonical\_\allowbreak measure} & the ``canonical discriminant measure'' claim & MARGINAL-WEIGHT (killed) & 26/26\\
T55 & \texttt{\scriptsize rtf\_\allowbreak period\_\allowbreak pairing} & canonical distributional limit of the finite RTF & DISTRIBUTIONAL-RTF-TYPED $\to$ \vref{v539} & 36/36\\
T56 & \texttt{\scriptsize rtf\_\allowbreak pairing\_\allowbreak sharpening} & all T55 identifications to sub-percent & SHARPENED & 27/27\\
T57 & \texttt{\scriptsize infinite\_\allowbreak rtf\_\allowbreak packing} & both infinity directions in one double series $Z(s,w)$ & PACKED & 24/24\\
T58 & \texttt{\scriptsize zeta\_\allowbreak rtf\_\allowbreak xi\_\allowbreak residue} & GL(1) shadows of $Z(s,w)$ & GL1-SHADOW-CLASSICAL (route A closed) & 24/24\\
T59 & \texttt{\scriptsize weil\_\allowbreak gns\_\allowbreak identification} & residual factor in closed form & TRANSFORM-REQUIRED & 19/19\\
T60 & \texttt{\scriptsize plancherel\_\allowbreak descent} & Sato--Tate moment contraction & PARTIAL & 21/21\\
T61 & \texttt{\scriptsize st\_\allowbreak isotype\_\allowbreak gl1\_\allowbreak core} & the trivial Sato--Tate isotype is a GL(1) core & GL1-CORE-EXACT $\to$ \vref{v539} & 24/24\\
T62 & \texttt{\scriptsize core\_\allowbreak isolation} & $d$-twist mixing vanishes fibrewise & CORE-ISOLATED $\to$ \vref{v539} & 24/24\\
T63 & \texttt{\scriptsize weil\_\allowbreak core\_\allowbreak completion} & Plancherel correction irreducibly non-automorphic & EXTRA-TERMS $\to$ \vref{v539} & 33/33\\
T64 & \texttt{\scriptsize rtf\_\allowbreak stabilization} & $\det_2$ absorbs obstruction 2; the minus stays & HALF-STABLE & 34/34\\
T65 & \texttt{\scriptsize pi\_\allowbreak digits\_\allowbreak prime\_\allowbreak distribution} & prime distribution in the digits of $\pi$ & PI-NULL & 16/16\\
T66 & \texttt{\scriptsize pi\_\allowbreak prime\_\allowbreak four\_\allowbreak level} & the $\pi$/prime question on four levels & FOUR-LEVEL-NULL & 23/23\\
T67 & \texttt{\scriptsize amplitude\_\allowbreak dirac\_\allowbreak sqrt} & Dirac square root of the amplitude plane & DIRAC-SQRT-EXACT $\to$ \vref{v540} & 27/27\\
T68 & \texttt{\scriptsize geometric\_\allowbreak polarization} & geometric polarisation $b=N_+-N_-$ & RESHUFFLE-ONLY $\to$ \vref{v540} & 34/34\\
T69 & \texttt{\scriptsize mixed\_\allowbreak channel\_\allowbreak full\_\allowbreak weight} & even-$k$ deletion is a determinant invariant & DELETION-UNIVERSAL $\to$ \vref{v540} & 33/33\\
T70 & \texttt{\scriptsize linear\_\allowbreak measure\_\allowbreak weil} & positive linear $\Theta$ carrier; pure plus combination & LINEAR-PLUS-COMBINATION $\to$ \vref{v540} & 29/29\\
T71 & \texttt{\scriptsize transport\_\allowbreak terrain} & FE of the linear family exactly anchored & TERRAIN-MAPPED $\to$ \vref{v540} & 40/40\\
T72 & \texttt{\scriptsize cone\_\allowbreak enlargement} & the cone library saturates at 5/24 & COMPLEMENTARY-PAIR $\to$ \vref{v540} & 34/34\\
T73 & \texttt{\scriptsize continuous\_\allowbreak cone} & per-direction hybrid cones in the continuum & HYBRID-GAINS & 28/28\\
T74 & \texttt{\scriptsize door\_\allowbreak a\_\allowbreak spectral\_\allowbreak polarization} & the spectral side is provably sign-blind & VACUUM-STRUCTURE & 32/32\\
T75 & \texttt{\scriptsize door\_\allowbreak b\_\allowbreak lambda\_\allowbreak transport} & the gap functional $\lambda^*$ has its own structure & LAMBDA-STRUCTURED & 28/28\\
T76 & \texttt{\scriptsize hybrid\_\allowbreak universality} & the recipe certifies 91/91 adversarial Weil directions & RECIPE-UNIVERSAL-ON-BATTERY & 22/22\\
T77 & \texttt{\scriptsize matching\_\allowbreak lemma\_\allowbreak structure} & the matching lemma reduces to a restricted divisor bound & LEMMA-CLASSICAL-SHAPED & 21/21\\
T78 & \texttt{\scriptsize lemma\_\allowbreak window\_\allowbreak proof} & matching lemma exact-integer on $[4,10^6]$ & WINDOW-PROVED $\to$ \vref{v541} & 25/25\\
T79 & \texttt{\scriptsize transport\_\allowbreak ledger} & the transport ledger closes; I5 named and typed & LEDGER-CLOSES $\to$ \vref{v541} & 22/22\\
T80 & \texttt{\scriptsize tail\_\allowbreak correlation\_\allowbreak lemma} & character-exact signed envelope; tail to $N^*\!\approx\!10^{23}$ & RESERVE-PARTIAL $\to$ \vref{v541} & 29/29\\
T81 & \texttt{\scriptsize recipe\_\allowbreak coherent\_\allowbreak avoidance} & reachability theorem: coherent targets need coherent $m$ & AVOIDANCE-FAILS $\to$ \vref{v541} & 31/31\\
T82 & \texttt{\scriptsize heat\_\allowbreak arch} & $\Darch$ is the internal $\Gamma$ difference & ARCH-INTERNAL $\to$ \vref{v541} & 22/22\\
T83 & \texttt{\scriptsize fe\_\allowbreak symmetrization} & symmetrisation absorbed; the wall is FE-transversal & INVARIANT-NULL & 27/27\\
T84 & \texttt{\scriptsize gaussian\_\allowbreak lift} & Gr\"ossencharacter phases give $L(1,\lambda_k)$ convergence & LIFT-WORKS-UNANCHORED & 29/29\\
T85 & \texttt{\scriptsize lambda\_\allowbreak equivariant\_\allowbreak design} & the $\lambda$ design closes the coherent class & LEMMA-CLOSES-LAMBDA $\to$ \vref{v541} & 24/24\\
T86 & \texttt{\scriptsize correlated\_\allowbreak cancellation} & the $Q$-pairing involution closes the non-coherent tail & LEMMA-FULLY-CLOSED & 29/29\\
T87 & \texttt{\scriptsize i5\_\allowbreak connes\_\allowbreak dictionary} & I5 one-family form vs.\ the semi-local Connes structure & DICTIONARY-EXACT-CORE & 22/22\\
T88 & \texttt{\scriptsize i5\_\allowbreak tightness} & eight near-vertical tight curves; $\Gamma$-side spacing & TIGHT-SET-PARAMETRIZED & 27/27\\
T89 & \texttt{\scriptsize sonin\_\allowbreak crossing} & window compression is Bombieri's classical object & CROSSING-MAPPED & 19/19\\
T90 & \texttt{\scriptsize residual\_\allowbreak dissection} & the residual vectors are explicit Gaussian modes & CORE-DISSECTED & 17/17\\
T91 & \texttt{\scriptsize thin\_\allowbreak band\_\allowbreak analytic} & the band is the one-atom zone; atom rescue & BAND-PARTIAL (3/4 closed) & 19/19\\
T92 & \texttt{\scriptsize zone\_\allowbreak extension\_\allowbreak proof} & the finite-block route is the wrong instrument & T-SKELETON & 36/36\\
T93 & \texttt{\scriptsize band\_\allowbreak reconciliation} & third independent implementation of the T89/T91 band map & MIXED (map recalibrated) & 41/41\\
T94 & \texttt{\scriptsize prime\_\allowbreak blind\_\allowbreak demo} & 753/753 primes of an unseen window from lattice counting & BLIND-100 & 16/16\\
T95 & \texttt{\scriptsize continuum\_\allowbreak extension} & C1 chain proved; band top behaves as a positivity edge & T-CONTINUUM-NUMERIC & 28/28\\
T96 & \texttt{\scriptsize relay\_\allowbreak test} & the T95 edge was a map artefact; the relay is real & EDGE-ARTIFACT $+$ RELAY-CONFIRMED & 21/21\\
T97 & \texttt{\scriptsize relay\_\allowbreak induction} & the induction step has proof shape; one cross-block lemma & ALIGNMENT-ONLY & 105/105\\
T98 & \texttt{\scriptsize cross\_\allowbreak lemma} & the $\sqrt{\lambda_0}$ law is Douglas range inclusion (circular) & LAW-CONFIRMED-MECHANISM-OPEN & 44/44\\
T99 & \texttt{\scriptsize dk\_\allowbreak recursion} & exact mirror-parity selection rule; recursive inequality & DECAY-LAW-FOUND & 23/23\\
T100 & \texttt{\scriptsize bulk\_\allowbreak remainder} & closure 11/24 $\to$ 24/24 samples; the drift was a grid artefact & REMAINDER-CLOSES-ZONES & 27/27\\
T101 & \texttt{\scriptsize k\_\allowbreak asymptotics} & 16 zones resolved; primitives flat, instrument loses the race & CROWDING-TRENDS & 31/31\\
T102 & \texttt{\scriptsize arithmetic\_\allowbreak bound} & the onset is manufactured by the Schur dressing; $C/g$ was a proxy & MECHANISM-IDENTIFIED & 42/42\\
T103 & \texttt{\scriptsize instrument} & slope halves; closure map $7/64\to44/64$ & INSTRUMENT-IMPROVED & 29/29\\
T104 & \texttt{\scriptsize schur\_\allowbreak profile\_\allowbreak bound} & arm B: naive margin route dead, threshold split closes 16/16 & CHAIN-PARTIAL & 21/21\\
T104 & \texttt{\scriptsize schur\_\allowbreak profile\_\allowbreak chain} & arm A: independent twin of the same contract & CHAIN-PARTIAL & 47/47\\
T105 & \texttt{\scriptsize bare\_\allowbreak wing\_\allowbreak bound} & bare bound certified in closed form; avoidance is a theorem & ONE-OF-TWO & 28/28\\
T106 & \texttt{\scriptsize friedrichs\_\allowbreak angle} & parity split of the Weil pole; hardness in the odd channel & DENSITY-MAPPED & 32/32\\
T107 & \texttt{\scriptsize odd\_\allowbreak channel\_\allowbreak closure} & Sherman--Morrison/Woodbury reduction to $r=\kappa/\varepsilon$ & SCALAR-TRACTABLE & 30/30\\
T108 & \texttt{\scriptsize ratio\_\allowbreak certificate} & $\varepsilon$ is an exact $Q|_{\mathrm{odd}}$ energy; (R) on two scalars; first non-vacuous margin chain & EPSILON-IDENTITY & 44/44\\
T109 & \texttt{\scriptsize boundary\_\allowbreak decay} & cancellation, not decay; both scalars certified (graded cap $+$ residual certificate); one strict-margin input left & BOUNDARY-CERTIFIED & 29/29\\
T110 & \texttt{\scriptsize margin\_\allowbreak propagation} & the induction circle closes on the measured zones (certified base case $+$ graded Loewner minorant; atom entry free); three gaps named & MARGIN-PROPAGATES & 28/28\\
T111 & \texttt{\scriptsize deep\_\allowbreak zone\_\allowbreak stress} & the crossing measured, not extrapolated ($n^*\!\approx\!462$, upper bound); three walls (margin / twin-prime ladder / requirement); mechanism never fails, $117/117$ & CROSSING-CONFIRMED & 23/23\\
T112 & \texttt{\scriptsize adaptive\_\allowbreak scaling} & gap-coupled frame: ladder wall and $\omega$ depth-dependence fall structurally; margin wall frame-invariant ($-0.974$); trade to [P1]/[P2], prime-gap inputs disclosed & SCALING-PARTIAL & 20/20\\
T113 & \texttt{\scriptsize limit\_\allowbreak operator} & the wall is currency-invariant ($-1.168$ in all five currencies) but measures the discretisation: no continuum gap, floor a $10^{-7}$--$10^{-4}$ cancellation, need109 an artifact margin; new balance [P1]--[P4] & SUBSTANCE-CONFIRMED & 27/27\\
T114 & \texttt{\scriptsize margin\_\allowbreak free\_\allowbreak step} & the margin-free step: Albert/pseudo-inverse Schur certifies $27/27$ steps ($11$ past the old wall, to $n=1331$) and all $7$ torn zones; the wall was an $O(1)$ numerator over an artifact floor; (R) grid-bound, deleted ([P5]); new two-object core & WALL-DISSOLVES & 22/22\\
T115 & \texttt{\scriptsize schur\_\allowbreak transport} & transport certifies only mild refinement (split at $\rho^*\!\approx\!1.8$; on nested ladders $\lambda_{\min}(S)$ itself falls like $\rho^{-1.7}$), but the two-scale compression breaks the cap: margin-free step at $n=155\,921$, ten-step chains, stopper always cost; remaining list three points, one inequality (Szeg\H{o}--Levinson) & TRANSPORT-BLOCKED & 26/26\\
\end{xltabular}
}

\section{Symbols used in Sections~\ref{sec:i5}--\ref{sec:certification}}\label{app:symbols}
\noindent All of these are sandbox objects; they are listed so the figures and formulas above can be
read without the diary.

{\small\noindent
\begin{tabularx}{\linewidth}{@{}>{\raggedright\arraybackslash}p{3.0cm} Y@{}}
\toprule
\textbf{Symbol} & \textbf{Meaning}\\
\midrule
$\Qweil$, $\Qcert$, $\Delta_2$, $\Darch$ & the four terms of the transport ledger;
$\Qweil=\Qcert+\Darch+\Delta_2$ closes at $4.4\times10^{-16}$ (\vref{v541})\\
$\Afam$, $\Ashift$ & the two archimedean halves of the heat family; $\Darch=\Afam-\Ashift$ (T82)\\
I5 & the one remaining coupled inequality, $\Qcert+\Delta_2+\Afam-\Ashift\ge0$ (T79)\\
$\mu_k$, $w_k$, $g_k$ & atom strength, handover window width, atom spacing at zone $k$\\
$D_k(\alpha)$ & the Schur-complement defect of the induction step; target $D_k\le\mu_k/2$ (T98)\\
$\sigma_k(\delta)$ & the $E_-$ Schur profile of the genuine Weil form just above atom entry (T102)\\
$E_-,E_0,E_+$ & the three atom eigenspaces; the $k$-th atom acts as $\mathrm{diag}(-\tfrac12,0,+\tfrac12)$\\
$\mathrm{bare}_k$, $b_0$ & $\lambda_{\min}(\Qfull|_{E_-})$ and its certified excess over $\mu_k/2$ (T105)\\
$L$, $\Lambda_-$ & the soft dressing scalar and its free cap ($5.63$--$7.52$, divergent) (T105)\\
$J$ & the mirror-parity involution; $JQJ=Q$, $JB=-BR$ (T105)\\
$\rho_\pm$ & the Friedrichs-angle cosines of the even / odd channel (T106)\\
(R) & $\Qfull(\alpha_k)|_{J=-1}\succeq(\mu_k/2)P_-|_{J=-1}$, the remaining Loewner statement (T106)\\
$s^*$, $\tau$, $\kappa$, $\varepsilon$ & Sherman--Morrison scalar, its Woodbury split
$s^*=\tau+\kappa$, and $\varepsilon=1-\tau$ (T107)\\
$r$ & the residual ratio $\kappa/\varepsilon$; (R) $\iff r\le1$; measured $0.0053$--$0.1814$ (T107)\\
$x$, $h$, $x_0$ & the pole response $x=T_{\mathrm{odd}}^{-1}\tilde t$, its demand compression
$h=V^{\mathsf T}x$, and the single edge entry that carries $\|h\|$ on $8$ zones (T108)\\
$\omega$, $\omega_{\mathrm{cert}}$ & the wing scalar $(\mu/2)\lambda_{\max}(\Gamma)$ with
$\Gamma=V^{\mathsf T}T_{\mathrm{odd}}^{-1}V$, and its graded-matrix-cap certificate (T108/T109)\\
$m_k$, $\mathrm{need109}$ & the strict odd-channel margin $\lambda_{\min}(Q|_{\mathrm{odd}})$ at zone
$k$, and the explicit demand $(\mu/2)H_{\mathrm{cert}}^2/((1-\omega_{\mathrm{cert}})\kappa)$ of the
T109 chain; $m_k\ge\mathrm{need109}$ on $16/16$ measured zones (T110)\\
\bottomrule
\end{tabularx}
}

\end{document}
